% 译自 arXiv:2306.13589v2 原 tex 第 2960–3572 行
\section{凝聚层的导出范畴}
\thispagestyle{plain}

\takeaway{%
  \textbf{物理图像：}本节为同调镜像对称的 B-侧搭好舞台：凝聚层的导出范畴在物理上就是 B-模型里 D-膜所构成的范畴，而层之间的 $\text{Ext}$ 群给出 D-膜之间的态射空间。\\
  \textbf{数学工具：}环空间上模层的导出范畴（内射消解与右导出函子）、作为张量函子适配类的平坦层、层上同调与 \v{C}ech 上同调及其在诺特分离概形上的比较。\\
  \textbf{主要结论：}模层范畴有足够多的内射对象，故 $\mathsf{R}\Gamma$、$\mathsf{R}f_*$、$\mathsf{R}\mathcal{H}om$ 等右导出函子存在；在诺特分离概形上、对拟凝聚层，\v{C}ech 上同调可用以计算层上同调；在光滑射影簇上，凝聚层有界导出范畴上的这些导出函子满足投影公式、伴随公式、平坦基变换与 Künneth 公式。%
}

\subsection{模层上的导出函子}\label{ch6.1}

让我们暂且退一步，重新把注意力放回模层上；只有在这一基础工作完成之后，我们才会在第 \ref{ch6.2} 节进行到一半时回到概形与凝聚层。

根据定理 \ref{Sh(X,Ab)isabelian}、注 \ref{sheafmodulesrem}a. 以及我们在第 2 章所作的分析，对取值于任意阿贝尔范畴 $\mathsf{A}$ 的、拓扑空间 $X$ 上的层，我们可以合法地谈论其导出范畴 $\mathsf{D}(\mathsf{Sh}(X;\mathsf{A}))$；相应地，对环空间 $(X,\mathcal{S})$ 上的 $\mathcal{S}$-模，也可以谈论其导出范畴 $\mathsf{D}(\mathsf{Sh}(X;{}_\mathcal{S}\mathsf{Mod}))$。我们已经知道若干个联系它们底层层范畴的函子，那么不妨看看它们是否诱导出导出函子。

\begin{Pro}\label{enoughinjectivemodules}
设 $(X,\mathcal{S})$ 为环空间。则 $\mathsf{Sh}(X;{}_\mathcal{S}\mathsf{Mod})$ 有足够多的内射对象 \textup(从而由注 \ref{sheafmodulesrem}b.，对任意阿贝尔范畴 $\mathsf{A}$，$\mathsf{Sh}(X;\mathsf{A})$ 也有足够多的内射对象\textup)。
\end{Pro}

\begin{proof}
根据定义 \ref{injprojobj}，需要证明的是：任意 $\mathcal{F}\in\text{obj}(\mathsf{Sh}(X;{}_\mathcal{S}\mathsf{Mod}))$ 都容许一个由内射 $\mathcal{S}$-模构成的右消解，特别地，$\mathcal{F}$ 单射入某个作为对象为内射的 $\mathcal{I}\in\text{obj}(\mathsf{Sh}(X;{}_\mathcal{S}\mathsf{Mod}))$。

首先注意到，对每个 $x\in X$，茎 $\mathcal{F}_x\in\text{obj}({}_{\mathcal{S}_x}\!\mathsf{Mod})$ 单射入另一个 $\mathcal{S}_x$-模 $I_x$，而后者作为 ${}_{\mathcal{S}_x}\!\mathsf{Mod}\equiv\mathsf{Sh}(\{x\};{}_\mathcal{S}\mathsf{Mod})$ 的对象是内射的\footnote{这是标准代数中的已知事实：给定环 $R$，每个 $R$-模都同构于某个内射 $R$-模的子模。}。记嵌入为 $\iota_x:\{x\}\hookrightarrow X$，于是 $(\iota_x)_*I_x\in\text{obj}(\mathsf{Sh}(X;{}_\mathcal{S}\mathsf{Mod}))$ 就是相应的摩天层（见例 \ref{skyscrapersheaf}）；今考虑层 $\mathcal{I}\coloneqq\prod_{x\in X}(\iota_x)_*I_x$（它是 $\mathcal{S}$-模，因为 $\mathsf{Sh}(X;{}_\mathcal{S}\mathsf{Mod})$ 是阿贝尔范畴，从而包含其对象的积！），其茎满足 $\mathcal{I}_x\cong I_x$。

给定任意另一个 $\mathcal{S}$-模 $\mathcal{G}$，我们有典范同构 $\text{Hom}_\mathcal{S}(\mathcal{G},\mathcal{I})\cong\prod_{x\in X}\text{Hom}_\mathcal{S}(\mathcal{G},(\iota_x)_*I_x)$。反过来，对每个 $x\in X$ 有 $\text{Hom}_\mathcal{S}(\mathcal{G},(\iota_x)_*I_x)\cong\text{Hom}_{\mathcal{S}_x}(\mathcal{G}_x,I_x)$，因为左端任意 $\mathcal{S}$-模态射 $\chi$ 在集合层面要么是到 $I_x$ 的常值映射 $\chi_U\cong\chi_x$（当 $x\in U$ 时），要么是零映射（否则）。取 $\mathcal{G}=\mathcal{F}$，所得的同构 $\text{Hom}_\mathcal{S}(\mathcal{F},\mathcal{I})\cong\prod_{x\in X}\text{Hom}_{\mathcal{S}_x}(\mathcal{F}_x,I_x)$ 告诉我们：局部单射 $\mathcal{F}_x\hookrightarrow I_x$ 诱导出一个 $\mathcal{S}$-模态射 $\mathcal{F}\rightarrow\mathcal{I}$，而它在茎上单射，故自身也是单射。

现在，回顾正合茎函子 $\square_x: \mathsf{Sh}(X;{}_\mathcal{S}\mathsf{Mod})\rightarrow{}_{\mathcal{S}_x}\!\mathsf{Mod}$ 的存在性，我们观察到按构造有
\[
\text{Hom}_\mathcal{S}(\square,\mathcal{I})=\prod_{x\in X}(\text{Hom}_{\mathcal{S}_x}(\square,I_x)\circ\square_x):\mathsf{Sh}(X;{}_\mathcal{S}\mathsf{Mod})^\text{opp}\rightarrow\mathsf{Ab}\,.
\] 
但如引理 \ref{Homforinjprojobj} 所提示的，每个左正合函子 $\text{Hom}_{\mathcal{S}_x}(\square,I_x):{}_{\mathcal{S}_x}\!\mathsf{Mod}^\text{opp}\rightarrow\mathsf{Ab}$ 实际上是正合的！因此，它与 $\square_x$ 的复合也正合，再取直积之后仍然正合。于是 $\text{Hom}_\mathcal{S}(\square,\mathcal{I})$ 正合。最后，再次由引理 \ref{Homforinjprojobj} 推知，$\mathcal{I}$ 必是内射的 $\mathcal{S}$-模层，此即所欲证者。
\end{proof}

因此，由引理 \ref{enoughinjprojobj} 与定义 \ref{derivedfunc}，我们可以考虑从 $\mathsf{Sh}(X;\mathsf{A})$ 或 $\mathsf{Sh}(X;{}_\mathcal{S}\mathsf{Mod})$ 出发的左正合函子的右导出函子：

\begin{Ex}
设 $f:(X,\mathcal{S}_X)\rightarrow (Y,\mathcal{S}_Y)$ 为环空间态射，并给定任意 $\mathcal{F}\in\text{obj}(\mathsf{Sh}(X;\mathsf{A}))$ 与 $\mathcal{G}\in\text{obj}(\mathsf{Sh}(X;{}_{\mathcal{S}_X}\!\mathsf{Mod}))$，则最突出的右导出函子的例子如下：
\begin{itemize}[leftmargin=0.5cm]
	\item (由例 \ref{exactfuncexample})
	\begin{equation}
		\begin{aligned}
			&\mathsf{R}\text{Hom}_{\mathsf{Sh}(X;\mathsf{A})}(\mathcal{F},\square):\mathsf{D^+}(\mathsf{Sh}(X;\mathsf{A}))\rightarrow\mathsf{D^+(Ab)}\,, \\ 
			&\mathsf{R}\text{Hom}_{\mathcal{S}_X}(\mathcal{G},\square):\mathsf{D^+}(\mathsf{Sh}(X;{}_{\mathcal{S}_X}\!\mathsf{Mod}))\rightarrow\mathsf{D^+(Ab)}\,. \\
		\end{aligned}
	\end{equation}
	
	\item (由定义 \ref{moresheafofmodules})
	\begin{equation}\label{localExtfunc}\begin{footnotesize}
			\mkern-12mu\mathcal{E}xt_{\mathcal{S}_X}\!(\mathcal{G},\square)\!\coloneqq\!\mathsf{R}\mathcal{H}om_{\mathcal{S}_X}\!(\mathcal{G},\square)\!:\!\mathsf{D^+}(\mathsf{Sh}(X;{}_{\mathcal{S}_X}\!\mathsf{Mod}))\!\rightarrow\!\mathsf{D^+}(\mathsf{Sh}(X;{}_{\mathcal{S}_X}\!\mathsf{Mod})), \\
		\end{footnotesize}
	\end{equation}
	类似地还可定义阿贝尔群层上的版本。（关于 $\mathcal{E}xt$ 的更多内容见下文定义 \ref{localExt}。）
	
	\item (由引理 \ref{inversepushforwardexact})
	\begin{equation}
		\begin{aligned}
			&\mathsf{R}f_*:\mathsf{D^+}(\mathsf{Sh}(X;\mathsf{A}))\rightarrow\mathsf{D^+}(\mathsf{Sh}(Y;\mathsf{A}))\,, \\ &\mathsf{R}f^{-1}:\mathsf{D^+}(\mathsf{Sh}(Y;\mathsf{A}))\rightarrow\mathsf{D^+}(\mathsf{Sh}(X;\mathsf{A}))\,.
		\end{aligned}
	\end{equation}
	
	\item (由推论 \ref{pullbackpushforwardexact})
	\begin{equation}\label{Rflowerstar}
		\mathsf{R}f_*:\mathsf{D^+}(\mathsf{Sh}(X;{}_{\mathcal{S}_X}\!\mathsf{Mod}))\rightarrow\mathsf{D^+}(\mathsf{Sh}(Y;{}_{\mathcal{S}_Y}\!\mathsf{Mod}))
	\end{equation} 
	
	\item (由引理 \ref{sectfuncexact}，对任意开集 $U\subset X$) 
	\begin{equation}\label{RGammaU}
		\mathsf{R}\Gamma(U,\square):\mathsf{D^+}(\mathsf{Sh}(X;\mathsf{A}))\rightarrow\mathsf{D^+(A)}\,.
	\end{equation}	  
\end{itemize}\vspace*{-0.2cm}
\hfill $\blacklozenge$ 
\end{Ex}

% [Tutor Note: 此处右导出函子存在的依据是命题 \ref{enoughinjectivemodules}：$\mathsf{Sh}(X;\mathsf{A})$ 有足够多的内射对象。注意这里 $U\subset X$ 不必是仿射的，$\mathsf{R}\Gamma(U,\square)$ 也只是左正合函子 $\Gamma(U,\square)$ 的右导出；若 $U$ 仿射而 $X$ 是分离概形，则由本节关于 \v{C}ech 上同调计算层上同调的定理可知它在高阶为零。]
\begin{Rem}\label{sectfuncispushforward}
以下几件事值得考虑，先给出一点警告。
\begin{itemize}[leftmargin=0.5cm]
	\item (重要警告) 我们\textit{不能}用同样的方法构造左正合函子 $\text{Hom}_{\mathsf{Sh}(X;\mathsf{A})}(\square,\mathcal{F})$、$\text{Hom}_{\mathcal{S}_X}(\square,\mathcal{G})$ 与 $\mathcal{H}om_{\mathcal{S}_X}\!(\square,\mathcal{G})$ 的右导出函子，因为这些是反变函子，也就是说我们要求源范畴有足够多的投射对象而非内射对象。然而这一点并不成立，我们稍后在考察左导出函子是否存在时会提到。
	
	\item 考虑上述右导出 Hom 函子的经典导出函子（见定义 \ref{classicalderivedfunc}），由命题 \ref{ExtRHom} 我们得到 \eqref{Extmodules} 处提到的层论 Ext 函子的另一种刻画（这一般被当作它们的默认定义）：
	\begin{equation}\label{ExtRHomformodules}
		\begin{aligned}
			&\text{Ext}_{\mathsf{Sh}(X;\mathsf{A})}^n(\mathcal{F},\square)\cong\mathsf{R}^n\text{Hom}_{\mathsf{Sh}(X;\mathsf{A})}(\mathcal{F},\square):\mathsf{Sh}(X;\mathsf{A})\rightarrow\mathsf{Ab}\,, \\ 
			&\text{Ext}_{\mathcal{S}_X}^n(\mathcal{G},\square)\cong\mathsf{R}^n\text{Hom}_{\mathcal{S}_X}(\mathcal{G},\square):\mathsf{Sh}(X;{}_{\mathcal{S}_X}\!\mathsf{Mod})\rightarrow\mathsf{Ab}\,, \\ 
		\end{aligned}  
	\end{equation}
	对每个 $n\in\mathbb{N}$ 成立（借助定理 \ref{Extprop}，我们知道当 $n<0$ 时这些是平凡函子，而当 $n=0$ 时则恢复原本的 Hom 函子）。
	
	\item 函子 \eqref{Rflowerstar} 对每个 $n\in\mathbb{Z}$ 诱导出经典的 $\mathsf{R}^nf_*:\mathsf{Sh}(X;{}_{\mathcal{S}_X}\!\mathsf{Mod})\rightarrow\mathsf{Sh}(Y;{}_{\mathcal{S}_Y}\!\mathsf{Mod})$。人们称它们为\textit{高阶直接像函子}；并且对任意 $\mathcal{S}_X$-模 $\mathcal{F}$，称 $\mathsf{R}^nf_*(\mathcal{F})=H^n(\mathsf{R}f_*(\mathcal{F}))\in\text{obj}(\mathsf{Sh}(Y;{}_{\mathcal{S}_Y}\!\mathsf{Mod}))$ 为 $\mathcal{F}$ 的\textbf{高阶直接像}\index{higher direct image}。
	
	\item 特别地取 $(Y,\mathcal{S}_Y)=(\{y\},\underline{\mathbb{Z}})$，所得的推前函子 $f_*:\mathsf{Sh}(X;{}_{\mathcal{S}_X}\!\mathsf{Mod})\rightarrow\mathsf{Sh}(\{y\};{}_{\underline{\mathbb{Z}}}\mathsf{Mod})\cong\mathsf{Sh}(\{y\};\mathsf{Ab})\equiv\mathsf{Ab}$ 不是别的，正是引理 \ref{sectfuncexact} 中的\textbf{整体截面函子}\index{functor!global sections} $\Gamma(X,\square):\mathsf{Sh}(X;{}_{\mathcal{S}_X}\!\mathsf{Mod})\rightarrow\mathsf{Ab}$ 在 $\mathcal{S}_X$-模上的特化，因而左正合，其相应的右导出函子为 $\mathsf{R}\Gamma(X,\square)':\mathsf{D^+}(\mathsf{Sh}(X;{}_{\mathcal{S}_X}\!\mathsf{Mod}))\rightarrow \mathsf{D^+}(\mathsf{Ab})$。\footnote{$\mathsf{R}\Gamma(X,\square)$ 旁的 $'$ 用以区别于 $\mathsf{D^+}(\mathsf{Sh}(X;\mathsf{A}))$ 上的右导出函子 $\mathsf{R}\Gamma(X,\square)$，后者由 \eqref{RGammaU} 中取 $U=X$ 得到。这只是一个临时记号，其动机见下文命题 \ref{sheafcohomologyisabelian}。}
\end{itemize} 
\end{Rem}

我们尤其强调 \eqref{localExtfunc} 中定义的右导出函子：

\begin{Def}\label{localExt}
设 $(X,\mathcal{S})$ 为环空间。对 $n\in\mathbb{Z}$ 及给定的 $\mathcal{F},\mathcal{G}\in\text{obj}(\mathsf{Sh}(X;{}_{\mathcal{S}}\mathsf{Mod}))$，$\mathcal{H}om_{\mathcal{S}}(\mathcal{F},\square)$ 的第 $n$ 个经典导出函子定义了\textbf{局部 $\mathcal{E}xt^n$-层}\index{sheaf!local $\mathcal{E}xt^n$} 
\begin{equation}
	\mathcal{E}xt_{\mathcal{S}}^n(\mathcal{F},\mathcal{G})\coloneqq H^n(\mathsf{R}\mathcal{H}om_{\mathcal{S}}(\mathcal{F},\mathcal{G}))\in\text{obj}(\mathsf{Sh}(X;{}_{\mathcal{S}}\mathsf{Mod}))\,,
\end{equation} 
它也可作为如下预层的层化而得到：该预层由对应 $U\subset X\text{ open}\mapsto \text{Ext}_{\mathcal{S}|_U}^n(\mathcal{F}|_U,\mathcal{G}|_U)$ 给出。

正如推论 \ref{whenclassicalderivedtrivial} 所预期的，$\mathcal{E}xt_{\mathcal{S}}^0(\mathcal{F},\mathcal{G})\cong \mathcal{H}om_{\mathcal{S}}(\mathcal{F},\mathcal{G})$，且当 $n<0$ 时\break $\mathcal{E}xt_{\mathcal{S}}^n(\mathcal{F},\mathcal{G})=0$；而 $\mathcal{E}xt_{\mathcal{S}}^0(\mathcal{S},\mathcal{G})\cong \mathcal{G} $，且当 $n\neq 0$ 时 $\mathcal{E}xt_{\mathcal{S}}^n(\mathcal{S},\mathcal{G})=0$（因为 $\mathcal{H}om_{\mathcal{S}}(\mathcal{S},\square)$ 是恒等函子；参看定义 \ref{moresheafofmodules}）。
\end{Def}

\begin{Rem}\label{localExtRem}
我们陈述关于局部 $\mathcal{E}xt$-层的两条性质（不加证明；建议读者查阅 \cite[section III.6]{[Har77]}）。
\begin{itemize}[leftmargin=0.5cm]
	\item 给定 $\mathcal{F},\mathcal{G}\in\text{obj}(\mathsf{Sh}(X;{}_{\mathcal{S}}\mathsf{Mod}))$，可以通过考虑 $\mathcal{F}$ 的局部自由投射消解来计算 $\mathcal{E}xt_{\mathcal{S}}^n(\mathcal{F},\mathcal{G})$，即考虑正合序列 $\mathcal{L}^\bullet\rightarrow\mathcal{F}\rightarrow 0$，其中每个 $\mathcal{L}^i\in\text{obj}(\mathsf{Sh}^\mathsf{lf}_{r_i}(X;{}_{\mathcal{S}}\mathsf{Mod}))$ 是秩为有限 $r_i<\infty$ 的局部自由层。于是 $\mathcal{E}xt_{\mathcal{S}}^n(\mathcal{F},\mathcal{G})\cong H^n(\mathcal{H}om_\mathcal{S}(\mathcal{L}^\bullet,\mathcal{G}))$。（当 $X$ 是 $\text{Spec}(R)$ 上的拟射影概形、$R$ 为诺特环时，这一点尤其有用，因为任意 $\mathcal{F}\in\text{obj}(\mathsf{Coh}(X;{}_{\mathcal{S}}\mathsf{Mod}))$ 都有这样的消解。）
	
	\item 给定 $\mathcal{F},\mathcal{G}\in\text{obj}(\mathsf{Sh}(X;{}_{\mathcal{S}}\mathsf{Mod}))$ 及某个 $\mathcal{L}\in\text{obj}(\mathsf{Sh}^\mathsf{lf}_r(X;{}_{\mathcal{S}}\mathsf{Mod}))$，有
	\begin{align*}
		&\text{Ext}^n(\mathcal{F}\otimes_\mathcal{S}\mathcal{L},\mathcal{G})\cong \text{Ext}^n(\mathcal{F},\mathcal{L}^\vee\otimes_\mathcal{S}\mathcal{G})\qquad\text{and} \\
		&\mathcal{E}xt_\mathcal{S}^n(\mathcal{F}\otimes_\mathcal{S}\mathcal{L},\mathcal{G})\cong \mathcal{E}xt_\mathcal{S}^n(\mathcal{F},\mathcal{L}^\vee\otimes_\mathcal{S}\mathcal{G})\cong \mathcal{E}xt_\mathcal{S}^n(\mathcal{F},\mathcal{G})\otimes_\mathcal{S}\mathcal{L}^\vee\,.
	\end{align*}     
\end{itemize}
\end{Rem}

% [Tutor Note: 局部 $\mathcal{E}xt$ 层与整体 $\text{Ext}$ 群切勿混淆：前者是 $\mathcal{H}om$ 的右导出函子取上同调层，取值于层范畴；后者是对整体截面取导出再取上同调，取值于 $\mathsf{Ab}$。二者由引理 \ref{sheafcohomthroughext} 中的关系 $H^n(X,\mathcal{F})\cong\textup{Ext}^n_{\mathcal{S}}(\mathcal{S},\mathcal{F})$ 联系，而不是相等。]
\vspace*{0.5cm}
\noindent 若要转而得到层上已知右正合函子的左导出函子，我们理想上需要命题 \ref{enoughinjectivemodules} 的投射对偶版本。然而在一般环空间上这不成立！（例如，利用投射性图表可以轻易证明：对任意单点空间 $\{x\}$（其在 $X$ 中有连通邻域），常值层 $\underline{\mathbb{Z}}$ 经零延拓所得的 $\iota_*\underline{\mathbb{Z}}_x\in\textup{obj}(\mathsf{Sh}(X;{}_{\underline{\mathbb{Z}}}\mathsf{Mod}))$ 不是某个投射层的商。）

于是我们转而直接为所考察的右正合函子寻找适配类。例 \ref{Torfunc} 为张量函子提供了很好的思路（参看 \cite[Proposition III.8.5]{[GM03]}）：

\begin{Lem}\label{flatsareadapted}
设 $(X,\mathcal{S})$ 为环空间，$\mathcal{M}\in\textup{obj}(\mathsf{Sh}(X;\mathsf{Mod}_\mathcal{S}))$，给出函子 $\mathcal{M}\otimes_\mathcal{S}\square:\mathsf{Sh}(X;{}_\mathcal{S}\mathsf{Mod})\rightarrow \mathsf{Sh}(X;\mathsf{Ab}),\,\mathcal{N}\mapsto\mathcal{M}\otimes_\mathcal{S}\mathcal{N}$，它是右正合的 \textup(正如例 \ref{exactfuncexample} 中环上模的情形\textup)。则全体平坦层构成的类是它的适配类，其中 $\mathcal{N}\in\textup{obj}(\mathsf{Sh}(X;{}_\mathcal{S}\mathsf{Mod}))$ 称为在 $\mathcal{S}$ 上\textup{平坦}（\textup{flat}），若 $\square\otimes_\mathcal{S}\mathcal{N}$ 正合。\footnote{于是注意到：$\mathcal{N}$ 平坦当且仅当对每个 $x\in X$，茎 $\mathcal{N}_x$ 是平坦 $\mathcal{S}_x$-模。}

对右正合函子 $\square\otimes_\mathcal{S}\mathcal{N}:\mathsf{Sh}(X;\mathsf{Mod}_\mathcal{S})\rightarrow \mathsf{Sh}(X;\mathsf{Ab})$ 的镜像陈述同样成立；此外我们指出，当 $\mathcal{S}$ 属于 $\mathsf{Sh}(X;\mathsf{CommRings})$ 时，两个函子的靶范畴都是 $\mathsf{Sh}(X;\mathsf{Mod}_\mathcal{S})$。
\end{Lem}

\begin{proof}
这里最重要的验证是：任意 $\mathcal{F}\in\textup{obj}(\mathsf{Sh}(X;{}_\mathcal{S}\mathsf{Mod}))$ 都是平坦层的商。若对任意开子集 $U\subset X$，用 $(\mathcal{S}_U)_x\coloneqq\mathcal{S}_x$（当 $x\in U$）与 $(\mathcal{S}_U)_x\coloneqq 0$（否则）来定义 $\mathcal{S}_U\in\text{obj}(\mathsf{Sh}(X;{}_\mathcal{S}\mathsf{Mod}))$，则我们得到一个平坦左 $\mathcal{S}$-模，且通过 $\psi\mapsto\psi_U(1_{\mathcal{S}(U)})$ 有 $\text{Hom}_\mathcal{S}(\mathcal{S}_U,\mathcal{F})\cong \mathcal{F}(U)$（参看同构 \eqref{Homglobalsectiso}）。

现在考虑一族开集 $\{U_i\}_{i\in I}$ 及相应的截面 $a_i\in\mathcal{F}(U_i)$，使得对任意 $x\in X$，茎 $\mathcal{F}_x$ 作为 $\mathcal{S}_x$-模由那些满足 $x\in U_i$ 的芽 $(a_i)_x=[(U_i,a_i)]$ 生成。那么由前述同构 $\Gamma(\bigoplus_{i\in I}U_i,\mathcal{F})\cong$ $\text{Hom}_\mathcal{S}(\bigoplus_{i\in I}\mathcal{S}_{U_i},\mathcal{F})$ 所确定的 $\mathcal{S}$-模态射 $\phi:\bigoplus_{i\in I}\mathcal{S}_{U_i}\rightarrow\mathcal{F}$ 显然满射，因为它在茎上满射。但平坦层的任意直和自身仍平坦，由此推出 $\mathcal{F}\cong (\bigoplus_{i\in I}\mathcal{S}_{U_i})/\ker(\phi)$，即所欲求的平坦层之商。

根据定义 \ref{adaptedclass}，还需证明 $\mathcal{M}\otimes_\mathcal{S}\square$ 把平坦层的零调复形映为零调的阿贝尔群层复形。这完全类似于环上模的情形：先证明任意左 $\mathcal{S}_x$-模的短正合三元组中，若后两个对象平坦则第一个也平坦（见 \cite[Exercise III.7.5]{[GM03]}），再对每个 $x\in X$ 在茎的层面应用这一结果。
\end{proof}

% [Tutor Note: 此处"平坦"是逐茎判定的：$\mathcal{N}$ 平坦当且仅当所有茎 $\mathcal{N}_x$ 是平坦 $\mathcal{S}_x$-模。证明的最后一步正是把环上的平坦性判据（短正合三元组中后两个平坦则第一个平坦）搬到茎上；注意这一步用到 $\square_x$ 正合，且"零调复形"的判定也只能在茎上做，不能直接对整体截面做。]
于是我们得到如下一组例子：

\begin{Ex}\label{derivedtensorfunc}
设 $(X,\mathcal{S})$ 为环空间，$\mathcal{M}\in\textup{obj}(\mathsf{Sh}(X;\mathsf{Mod}_\mathcal{S}))$，相应地 $\mathcal{N}\in\textup{obj}(\mathsf{Sh}(X;{}_\mathcal{S}\mathsf{Mod}))$。则由引理 \ref{flatsareadapted} 与定义 \ref{derivedfunc}，存在左导出函子
\begin{equation}
	\begin{aligned}
		&\mathcal{M}\overset{\mathsf{L}}{\otimes}_\mathcal{S}\square\coloneqq\mathsf{L}(\mathcal{M}\otimes_\mathcal{S}\square):\mathsf{D^-}(\mathsf{Sh}(X;{}_\mathcal{S}\mathsf{Mod}))\rightarrow\mathsf{D^-}(\mathsf{Sh}(X;\mathsf{Ab}))\,, \\ &\square\overset{\mathsf{L}}{\otimes}_\mathcal{S}\mathcal{N}\coloneqq\mathsf{L}(\square\otimes_\mathcal{S}\mathcal{N}):\mathsf{D^-}(\mathsf{Sh}(X;\mathsf{Mod}_\mathcal{S}))\rightarrow\mathsf{D^-}(\mathsf{Sh}(X;\mathsf{Ab}))\,,
	\end{aligned}
\end{equation}
其经典对应物（见定义 \ref{classicalderivedfunc}）给出上同调层，记作 $\text{Tor}_n^\mathcal{S}(\mathcal{M},\mathcal{N})\coloneqq H^{-n}(\mathcal{M}\overset{\mathsf{L}}{\otimes}_\mathcal{S}\mathcal{N})\in\text{obj}(\mathsf{Sh}(X;\mathsf{Ab}))$（两个函子给出同样的层！），其中 $n\geq 0$，这与例 \ref{Torfunc} 的记号一致。

此外，给定环空间态射 $f:(X,\mathcal{S}_X)\rightarrow (Y,\mathcal{S}_Y)$，其中 $\mathcal{S}_X,\mathcal{S}_Y\in\text{obj}(\mathsf{Sh}(X;\mathsf{CommRings}))$，由推论 \ref{pullbackpushforwardexact} 我们知道拉回函子是右正合的，于是得到左导出函子
\begin{equation}
	\mathsf{L}f^*:\mathsf{D^-}(\mathsf{Sh}(Y;{}_{\mathcal{S}_Y}\!\mathsf{Mod}))\rightarrow\mathsf{D^-}(\mathsf{Sh}(X;{}_{\mathcal{S}_X}\!\mathsf{Mod}))
\end{equation}
及其相应的层 $\mathsf{L}^{-n}f^*(\mathcal{G})=H^{-n}(\mathcal{S}_X\overset{\mathsf{L}}{\otimes}_{f^{\!-\!1}\mathcal{S}_Y} f^{-1}\mathcal{G})$，$n\geq 0$。\hfill $\blacklozenge$
\end{Ex}

% [Tutor Note: 注意 $\text{Tor}$ 的指标符号约定：作者用 $H^{-n}$ 而非 $H^n$ 定义 $\text{Tor}_n$，因为复形是下链的（微分降次）。同理 $\mathsf{L}^{-n}f^*$：拉回的左导出函子在 $\mathsf{D^-}$ 中把上同调次数整体下移。若换成同调编号，这些指标全部要加负号。]

\subsection{层上同调与 \v{C}ech 上同调}\label{ch6.2}

把定义 \ref{classicalderivedfunc} 应用于引理 \ref{sectfuncexact} 中的左正合整体截面函子，我们就得到\textit{层上同调}这一基本概念：

\begin{Def}\label{sheafcohomology}
设 $X$ 为拓扑空间，$\Gamma(X,\square):\mathsf{Sh}(X;\mathsf{Ab})\rightarrow\mathsf{Ab}$ 为整体截面函子，给出 $\mathsf{R}\Gamma(X,\square):\mathsf{D^+}(\mathsf{Sh}(X;\mathsf{Ab}))\rightarrow \mathsf{D^+}(\mathsf{Ab})$。给定 $n\in\mathbb{Z}$，我们称相应的经典第 $n$ 个右导出函子 $H^n(X,\square)\coloneqq \mathsf{R}^n\Gamma(X,\square)= H^n(\mathsf{R}\Gamma(X,\square)):\mathsf{Sh}(X;\mathsf{Ab})\rightarrow\mathsf{Ab}$ 为\textbf{第 $n$ 个层上同调函子}\index{functor!sheaf cohomology}，并称
\begin{equation}
	H^n(X,\mathcal{F})=H^n(\mathsf{R}\Gamma(X,\square))(\mathcal{F})\in\text{obj}(\mathsf{Ab})
\end{equation}
为 $\mathcal{F}\in\text{obj}(\mathsf{Sh}(X;\mathsf{Ab}))$ 的\textbf{第 $n$ 个层上同调群}\index{sheaf cohomology group}。特别地，由推论 \ref{whenclassicalderivedtrivial}，有 $H^0(X,\mathcal{F})\cong\Gamma(X,\mathcal{F})=\mathcal{F}(X)$，且当 $n<0$ 时 $H^n(X,\mathcal{F})\cong\{0\}$。
\end{Def}

\begin{Pro}\label{sheafcohomologyisabelian}
设 $(X,\mathcal{S})$ 为环空间，给出注 \ref{sectfuncispushforward} 末尾构造的函子 $\mathsf{R}\Gamma(X,\square)':\mathsf{D^+}(\mathsf{Sh}(X;{}_{\mathcal{S}}\mathsf{Mod}))\rightarrow \mathsf{D^+}(\mathsf{Ab})$。记 $\Phi:\mathsf{Sh}(X;{}_{\mathcal{S}}\mathsf{Mod})\rightarrow\mathsf{Sh}(X;\mathsf{Ab})$ 为把 $\mathcal{S}$-模层解释为纯粹的 $\underline{\mathbb{Z}}$-模层（从而是阿贝尔群层）的\textup(显然正合的\textup)遗忘函子，则有
\[
\mathsf{R}\Gamma(X,\square)'\cong\mathsf{R}\Gamma(X,\square)\circ\mathsf{R}\Phi:\mathsf{D^+}(\mathsf{Sh}(X;{}_{\mathcal{S}}\mathsf{Mod}))\rightarrow\mathsf{D^+}(\mathsf{Ab})\,,
\]
其中右端的 $\mathsf{R}\Gamma(X,\square)$ 定义在 $\mathsf{D^+}(\mathsf{Sh}(X;\mathsf{Ab}))$ 上，如 \eqref{RGammaU} 中取 $U=X$ 所得。

换言之，在计算层上同调时，我们可以忘掉拓扑空间 $X$ 及其层上的附加结构，而只关注定义 \ref{sheafcohomology} 的框架。\textup(于是在记号上，我们总可以不加撇号地写 $\mathsf{R}\Gamma(X,\square)$。\textup)   
\end{Pro}

\begin{proof}
显然 $\Gamma(X,\square)'=\Gamma(X,\square)\circ\Phi:\mathsf{Sh}(X;{}_{\mathcal{S}}\mathsf{Mod})\rightarrow\mathsf{Ab}$，故 $\mathsf{R}\Gamma(X,\square)'$ $=\mathsf{R}(\Gamma(X,\square)\circ\Phi):\mathsf{D^+}(\mathsf{Sh}(X;{}_{\mathcal{S}}\mathsf{Mod}))\rightarrow\mathsf{D^+}(\mathsf{Ab})$。我们的目标是应用命题 \ref{derivedfunccompo}，它给出典范同构 $\mathsf{R}(\Gamma(X,\square)\circ\Phi)\cong\mathsf{R}\Gamma(X,\square)\circ\mathsf{R}\Phi$。为此，我们需要两个类 $\mathcal{R}_1\subset\text{obj}(\mathsf{Sh}(X;{}_{\mathcal{S}}\mathsf{Mod}))$、$\mathcal{R}_2\subset\text{obj}(\mathsf{Sh}(X;\mathsf{Ab}))$，分别适配于 $\Phi$ 与 $\Gamma(X,\square)$，且满足 $\Phi(\mathcal{R}_1)\subset\mathcal{R}_2$。

第一个可行的候选是 $\mathcal{R}_1\coloneqq\{\mathcal{I}\in\text{obj}(\mathsf{Sh}(X;{}_{\mathcal{S}}\mathsf{Mod}))\mid \mathcal{I}\text{ is injective}\}$，依据是命题 \ref{enoughinjectivemodules} 与引理 \ref{enoughinjprojobj}。令 $\mathcal{R}_2\coloneqq\{\mathcal{F}\in\text{obj}(\mathsf{Sh}(X;\mathsf{Ab}))\mid \mathcal{F}\text{ is flabby}\}$，其中阿贝尔群层 $\mathcal{F}$ 称为「松弛的」（flabby），若对一切开集 $U\subset X$，限制映射 $\rho^X_U:\mathcal{F}(X)=\Gamma(X,\mathcal{F})\rightarrow\mathcal{F}(U)=\Gamma(U,\mathcal{F})$ 满射。给定层 $\mathcal{F}$，对应 $\mathcal{F}_\text{flabby}(U)\coloneqq\prod_{x\in U}\mathcal{F}_x$ 显然使 $\mathcal{F}_\text{flabby}$ 成为松弛层。设 $\mathcal{F}\in\mathcal{R}_1$，则 $\Phi(\mathcal{F})_\text{flabby}$ 按构造以 $\Phi(\mathcal{F})$ 为直和项，从而必为松弛层。这就验证了 $\Phi(\mathcal{R}_1)\subset\mathcal{R}_2$。

我们还须证明 $\mathcal{R}_2$ 适配于左正合函子 $\Gamma(X,\square)$。定义 \ref{adaptedclass} 的第一项要求显然满足：任意 $\mathcal{F}\in\text{obj}(\mathsf{Sh}(X;\mathsf{Ab}))$ 都是其 $\mathcal{F}_\text{flabby}\in\mathcal{R}_2$ 的子对象。保持零调性则需要更多工作。第一步是证明 $\mathsf{Sh}(X;\mathsf{Ab})$ 中任意短正合序列 $0\rightarrow\mathcal{F}\xrightarrow{\chi}\mathcal{G}\xrightarrow{\psi}\mathcal{H}\rightarrow 0$（其中 $\mathcal{F}\in\mathcal{R}_2$）都诱导出阿贝尔群的短正合序列 $0\rightarrow\Gamma(X,\mathcal{F})\xrightarrow{\Gamma(X,\chi)}\Gamma(X,\mathcal{G})\xrightarrow{\Gamma(X,\psi)}\Gamma(X,\mathcal{H})\rightarrow 0$，于是由整体截面函子的左正合性，只需验证 $\Gamma(X,\psi)$ 的满射性。这与其说困难不如说繁琐，故我们宁可请读者参阅 \cite[Theorem III.8.3a)]{[GM03]}。类似的方法可用以证明：若进一步 $\mathcal{G}\in\mathcal{R}_2$，则也有 $\mathcal{H}\in\mathcal{R}_2$。

今设 $\mathcal{F}^\bullet=(0\rightarrow\mathcal{F}^0\rightarrow\mathcal{F}^1\rightarrow...)\in\text{obj}(\mathsf{Kom^+}(\mathcal{R}_2))$，微分为 $d^\bullet$，且是零调的。这意味着 $\mathcal{Z}^n\coloneqq\ker(d^n)=\text{im}(d^{n-1})$，从而每个序列 $0\rightarrow\mathcal{Z}^n\hookrightarrow\mathcal{F}^n\xrightarrow{d^n}\mkern-20mu\rightarrow\mathcal{Z}^{n+1}\rightarrow 0$ 正合。由于 $\mathcal{Z}^0=0\in\mathcal{R}_2$ 且 $\mathcal{F}^0\in\mathcal{R}_2$，由刚才提到的事实也有 $\mathcal{Z}^1\in\mathcal{R}_2$。于是归纳可得 $\mathcal{Z}^n\in\mathcal{R}_2$ 对一切 $n\in\mathbb{N}$ 成立，进而由上述结论知 $0\rightarrow\Gamma(X,\mathcal{Z}^n)\hookrightarrow\Gamma(X,\mathcal{F}^n)\xrightarrow{\Gamma(X,d^n)}\mkern-20mu\rightarrow\Gamma(X,\mathcal{Z}^{n+1})\rightarrow 0$ 恒正合，故 $\text{im}(\Gamma(X,d^{n-1}))=\Gamma(X,\mathcal{Z}^n)=\ker(\Gamma(X,d^n))$。因此，对一切 $n\in\mathbb{N}$ 有 $H^n(\Gamma(X,\mathcal{F}^\bullet))=\ker(\Gamma(X,d^n))/\text{im}(\Gamma(X,d^{n-1}))\cong\{0\}$，即 $\Gamma(X,\mathcal{F}^\bullet)$ 零调。故 $\mathcal{R}_2$ 适配于 $\Gamma(X,\square)$。    
\end{proof}

% [Tutor Note: 这里的"松弛层"（flabby sheaf，也译作"软层"）与后文的 flasque 是同一概念，作者在证明中用 flabby、在引理中改用 flasque。判断一个层类是否"适配"（adapted）需两条：①每个对象都能嵌入该类中的对象；②该类中的零调复形经函子后仍零调。本证明正是逐条验证这两点，其中"松弛层的商仍松弛"这一关键步骤作者直接引用 \cite[Theorem III.8.3a)]{[GM03]} 而略去了图追踪。]
\begin{Lem}\label{sheafcohomthroughext}
设 $(X,\mathcal{S})$ 为环空间，$\mathcal{F}\in\textup{obj}(\mathsf{Sh}(X;{}_{\mathcal{S}}\mathsf{Mod}))$。则
\[
H^n(X,\mathcal{F})\cong\textup{Ext}_{\mathcal{S}}^n(\mathcal{S},\mathcal{F})\quad\textup{\big(}\cong\textup{Ext}_{\mathsf{Sh}(X;\mathsf{Ab})}^n(\underline{\mathbb{Z}},\mathcal{F}),\;\text{by Proposition \ref{sheafcohomologyisabelian}}\textup{\big)}
\]
对一切 $n\in\mathbb{Z}$ 成立。
\end{Lem}

\begin{proof}
这很直接：由 \eqref{ExtRHomformodules} 有 $\text{Ext}_{\mathcal{S}}^n(\mathcal{S},\mathcal{F})\cong\mathsf{R}^n\text{Hom}_{\mathcal{S}}(\mathcal{S},\mathcal{F})$，而由同构 \eqref{Homglobalsectiso} 与层上同调的定义有 $\mathsf{R}^n\text{Hom}_{\mathcal{S}}(\mathcal{S},\mathcal{F})\cong \mathsf{R}^n\Gamma(X,\mathcal{F})=H^n(X,\mathcal{F})$。（作为推论，注意到定理 \ref{Extprop} 中关于 $n\leq 0$ 的已知恒等式在此复核了平凡层上同调群的那些恒等式。）
\end{proof}

由引理 \ref{sheafcohomthroughext} 与定理 \ref{Extprop} 立即得到：

\begin{Cor}\label{longexactforsheafcohomology}
设 $(X,\mathcal{S})$ 为环空间。则 $\mathsf{Sh}(X;{}_{\mathcal{S}}\mathsf{Mod})$ 中任意短正合序列 $0\rightarrow\mathcal{F}_1\rightarrow\mathcal{F}_2\rightarrow\mathcal{F}_3\rightarrow 0$ 都诱导出层上同调群的长正合序列
\[
...\rightarrow H^n(X,\mathcal{F}_1)\rightarrow H^n(X,\mathcal{F}_2)\rightarrow H^n(X,\mathcal{F}_3)\rightarrow H^{n+1}(X,\mathcal{F}_1)\rightarrow...
\]
\end{Cor}

\vspace*{0.5cm}
\noindent 尽管有这些令人鼓舞的结果，层上同调在实践中仍相当难以计算。处理它的一个有用工具是\textit{\v{C}ech 上同调}（参见 \cite[section III.4]{[Har77]}）。

\begin{Def}
设 $X$ 为拓扑空间，带有开覆盖 $\mathfrak{U}=(U_i)_{i\in I}$，其中我们为 $I$ 选定一个良序。给定任意 $\mathcal{F}\in\text{obj}(\mathsf{Sh}(X;\mathsf{Ab}))$，我们可以定义阿贝尔群复形 $\check{C}^\bullet(\mathfrak{U},\mathcal{F})$：令
\begin{equation}\label{Cechcomplex}
	\check{C}^n(\mathfrak{U},\mathcal{F})\coloneqq\mkern-18mu\prod_{\begin{small}1\leq i_0<...<i_n\leq |I|\end{small}}\mkern-18mu\mathcal{F}(U_{i_0}\cap ...\cap U_{i_n})\in\text{obj}(\mathsf{Ab})\,,
\end{equation}
对 $n\geq 0$ 成立，并令 $n<0$ 时 $\check{C}^n(\mathfrak{U},\mathcal{F})\coloneqq 0$。于是任意 $s\in \check{C}^n(\mathfrak{U},\mathcal{F})$ 都具有形式 $s=(s_{i_0...i_n})$，其中 $s_{i_0...i_n}\in\mathcal{F}(U_{i_0}\cap ...\cap U_{i_n})$，对 $I$ 中每个允许的 $(n+1)$-元组 $1\leq i_0<...<i_n\leq |I|$ 成立。注意，若施加反对称性，我们也可以考虑一般的元组（可能无序且指标重复）：$s_{i_0...i_n}=(-1)^\sigma s_{\sigma(i_0)...\sigma(i_n)}$，其中 $\sigma\in S_{n+1}$ 满足 $1\leq \sigma(i_0)<...<\sigma(i_n)\leq |I|$（且指标重复时 $s_{i_0...i_n}=0$）。

微分 $\check{d}^\bullet:\check{C}^\bullet(\mathfrak{U},\mathcal{F})\rightarrow \check{C}^{\bullet+1}(\mathfrak{U},\mathcal{F})$ 由下式给出
\begin{equation}\label{Cechdifferential}
	(\check{d}^n s)_{i_0...i_{n+1}}\coloneqq \sum_{k=0}^{n+1}(-1)^k s_{i_0...\hat{i_k}...i_{n+1}}|_{U_{i_0}\cap...\cap U_{i_{n+1}}}\in \check{C}^{n+1}(\mathfrak{U},\mathcal{F})\,,
\end{equation}
其中 $s\in \check{C}^n(\mathfrak{U},\mathcal{F})$。（回想指标上的帽子表示省略该项，这就解释了为什么最后必须限制到 $U_{i_0}\cap...\cap U_{i_{n+1}}$。）容易验证 $\check{d}^{n+1}\circ \check{d}^n=0$，于是我们确实得到一个下方有界复形 $(\check{C}^\bullet(\mathfrak{U},\mathcal{F}),\check{d}^\bullet)\in\text{obj}(\mathsf{Kom^+(Ab)})$，即\textbf{\v{C}ech 复形}\index{cech@\v{C}ech complex}。
\end{Def}

\begin{Def}
设 $X$ 为带有开覆盖 $\mathfrak{U}=(U_i)_{i\in I}$ 的拓扑空间。则 $\mathcal{F}\in\text{obj}(\mathsf{Sh}(X;\mathsf{Ab}))$ 的\textbf{第 $n$ 个 \v{C}ech 上同调群}\index{cech@\v{C}ech cohomology group}定义为
\begin{equation}
	\check{H}^n(\mathfrak{U},\mathcal{F})\coloneqq H^n(\check{C}^\bullet(\mathfrak{U},\mathcal{F}))=\ker(\check{d}^n)/\text{im}(\check{d}^{n-1})\,.
\end{equation}
\end{Def}

% [Tutor Note: 这里出现了两套极易混淆的记号：$\check{C}^\bullet(\mathfrak{U},\mathcal{F})$ 是阿贝尔群复形、$\check{H}^n$ 是其上同调群；而花体的 $\check{\mathcal{C}}^\bullet(\mathfrak{U},\mathcal{F})$ 是层的复形，取整体截面后才回到 $\check{C}^\bullet$。读时务必留神花体与直体之别；下文引理 \ref{Cechinjres} 用的正是花体那套。]
\begin{Lem}\label{CechH0isglobalsect}
设 $X$ 为带有开覆盖 $\mathfrak{U}=(U_i)_{i\in I}$ 的拓扑空间，$\mathcal{F}\in\textup{obj}(\mathsf{Sh}(X;\mathsf{Ab}))$。则 $\check{H}^0(\mathfrak{U},\mathcal{F})\cong\mathcal{F}(X)$。
\end{Lem}

\begin{proof}
按定义，$\check{H}^0(\mathfrak{U},\mathcal{F})=\ker(\check{d}^0)/\{0\}\cong\ker(\check{d}^0\!:\check{C}^0(\mathfrak{U},\mathcal{F})\rightarrow \check{C}^1(\mathfrak{U},\mathcal{F}))$。任意 $s\in \check{C}^0(\mathfrak{U},\mathcal{F})=\prod_{i\in I}\mathcal{F}(U_i)$ 是一族截面 $s_i\in\mathcal{F}(U_i)$，且对 $i<j$ 有 $(\check{d}^0s)_{i,j}=(s_j-s_i)|_{U_i\cap U_j}$，于是 $\check{d}^0s\overset{!}{=}0$ 意味着每对截面在其交叠上一致。于是由层性质，存在唯一 $s\in\mathcal{F}(X)$ 限制为这些截面，故 $\ker(\check{d}^0)\subset \mathcal{F}(X)$；反向的包含由逆转上述论证得到。
\end{proof}

现在我们把 \v{C}ech 复形提升为层的复形。

\begin{Def}
设 $X$ 为带有开覆盖 $\mathfrak{U}=(U_i)_{i\in I}$ 的拓扑空间，$\mathcal{F}\in\text{obj}(\mathsf{Sh}(X;\mathsf{Ab}))$。$X$ 上的\textbf{层的 \v{C}ech 复形}\index{cech@\v{C}ech complex!of sheaves} $\check{\mathcal{C}}^\bullet(\mathfrak{U},\mathcal{F})$ 由下式给出
\begin{equation}
	\begin{aligned}
		&\check{\mathcal{C}}^n(\mathfrak{U},\mathcal{F})\coloneqq\mkern-18mu\prod_{\begin{small}1\leq i_0<...<i_n\leq |I|\end{small}}\mkern-18mu (\iota_{i_0...i_n})_*(\mathcal{F}|_{U_{i_0}\cap...\cap U_{i_n}})\in\text{obj}(\mathsf{Sh}(X;\mathsf{Ab}))\,, \\
		&\check{\mathcal{C}}^n(\mathfrak{U},\mathcal{F})(V)= \mkern-18mu\prod_{\begin{small}1\leq i_0<...<i_n\leq |I|\end{small}}\mkern-18mu \mathcal{F}|_{U_{i_0}\cap...\cap U_{i_n}}(\iota_{i_0...i_n}^{-1}(V))=\check{C}^n(\mathfrak{U}\cap V,\mathcal{F}|_V)\,,
	\end{aligned}
\end{equation}
其中 $n\geq 0$，$\iota_{i_0...i_n}:U_{i_0}\cap...\cap U_{i_n}\hookrightarrow X$ 为嵌入，给出推前 $(\iota_{i_0...i_n})_*:\mathsf{Sh}(U_{i_0}\cap...\cap U_{i_n};\mathsf{Ab})\rightarrow\mathsf{Sh}(X;\mathsf{Ab})$，而 $\mathfrak{U}\cap V=(U_i\cap V)_{i\in I}$ 是任意开子集 $V\subset X$ 的开覆盖。\footnote{回想第 $(i_0,...,i_n)$ 个直和项 $\mathcal{F}|_{U_{i_0}\cap...\cap U_{i_n}}(\iota_{i_0...i_n}^{-1}(V))$ 非零当且仅当 $V\subset U_{i_0}\cap...\cap U_{i_n}$。}

限制映射取显然的那些（同时也约束覆盖），而微分 $\check{d}^\bullet:\check{\mathcal{C}}^\bullet(\mathfrak{U},\mathcal{F})\rightarrow \check{\mathcal{C}}^{\bullet+1}(\mathfrak{U},\mathcal{F})$ 是全体映射 $\check{d}_V^n:\check{\mathcal{C}}^n(\mathfrak{U},\mathcal{F})(V)\rightarrow \check{\mathcal{C}}^{n+1}(\mathfrak{U},\mathcal{F})(V)$ 的汇集，其定义与 \eqref{Cechdifferential} 类似。于是我们得到真正的 $\check{\mathcal{C}}^\bullet(\mathfrak{U},\mathcal{F})\in\text{obj}\big(\mathsf{Kom^+}(\mathsf{Sh}(X;\mathsf{Ab}))\big)$。特别地，$\check{\mathcal{C}}^n(\mathfrak{U},\mathcal{F})(X)=\check{C}^n(\mathfrak{U},\mathcal{F})$。  
\end{Def}

\begin{Lem}\label{Cechinjres}
设 $X$ 为带有开覆盖 $\mathfrak{U}=(U_i)_{i\in I}$ 的拓扑空间。对任意 $\mathcal{F}\in\textup{obj}(\mathsf{Sh}(X;\mathsf{Ab}))$ 存在自然映射 $\varepsilon:\mathcal{F}\rightarrow \check{\mathcal{C}}^0(\mathfrak{U},\mathcal{F})$，使得层序列 $0\rightarrow\mathcal{F}\xrightarrow{\varepsilon}\check{\mathcal{C}}^0(\mathfrak{U},\mathcal{F})\xrightarrow{\check{d}^0}\check{\mathcal{C}}^1(\mathfrak{U},\mathcal{F})\rightarrow...$ 正合。也就是说，$\mathcal{F}$ 容许一个内射消解 $\mathcal{F}\xrightarrow{\varepsilon}\check{\mathcal{C}}^\bullet(\mathfrak{U},\mathcal{F})$。
\end{Lem}

\begin{proof}
由于 $\check{\mathcal{C}}^0(\mathfrak{U},\mathcal{F})=\prod_{i\in I} (\iota_i)_*(\mathcal{F}|_{U_i})$，我们令 $\varepsilon$ 为全体映射 $\mathcal{F}\rightarrow (\iota_i)_*(\mathcal{F}|_{U_i})$（$i\in I$）的直积，这使 $\varepsilon$ 成为自然的。于是每个（非平凡的）群同态 $\varepsilon_V:\mathcal{F}(V)\rightarrow \check{\mathcal{C}}^0(\mathfrak{U},\mathcal{F})(V)=\prod_{i\in I} \mathcal{F}|_V(U_i\cap V)$ 显然由层性质单射（层性质保证在 $V$ 上存在\textit{唯一}的整体截面，任意两个相等的相容子截面直积都必须粘合到它上面）。

对 $n\geq 1$，层 \v{C}ech 复形的正合性最好在茎上验证（这是等价的验证方式）。给定 $x\in X$，固定某个指标 $j\in I$ 使 $x\in U_j$，并在 $\check{\mathcal{C}}^\bullet(\mathfrak{U},\mathcal{F})_x$ 上定义同伦 $k^\bullet$，它由群同态 $k^n:\check{\mathcal{C}}^n(\mathfrak{U},\mathcal{F})_x\rightarrow \check{\mathcal{C}}^{n-1}(\mathfrak{U},\mathcal{F})_x,\, s_x=[(s,V)]\mapsto k^n(s_x)=[(k^ns,V)]$ 组成，其中 $(k^ns)_{i_0...i_{n-1}}\coloneqq s_{ji_0...i_{n-1}}$；这是合理的，因为 $x$ 的邻域 $V\subset X$ 取为满足 $V\subset U_j$（从而 $V\cap U_{i_0}\cap...\cap U_{i_{n-1}}=V\cap U_j\cap U_{i_0}\cap...\cap U_{i_{n-1}}$）。于是：
\begingroup
\allowdisplaybreaks
\begin{small}\begin{align*}
		(\check{d}&^{n-1}\circ k^n + k^{n+1}\circ \check{d}^n)(s)_{i_0...i_n} = \check{d}^{n-1}(k^ns)_{i_0...i_n} + k^{n+1}(\check{d}^ns)_{i_0...i_n} \\
		&= \sum_{k=0}^n(-1)^k(k^ns)_{i_0...\hat{i_k}...i_n}|_{U_{i_0}\cap...\cap U_{i_n}} + (\check{d}^ns)_{ji_0...i_n} \\
		&= \sum_{k=0}^n(-1)^k s_{ji_0...\hat{i_k}...i_n}|_{U_{i_0}\cap...\cap U_{i_n}} \!+\! \Big(\sum_{l=0}^{n+1}(-1)^l s_{i_0...\hat{i_l}...i_{n+1}}|_{U_{i_0}\cap...\cap U_{i_{n+1}}}\Big)_{\!j...i_n} \\
		&= \sum_{k=0}^n(-1)^k s_{ji_0...\hat{i_k}...i_n}|_{U_{i_0}\cap...\cap U_{i_n}} \!+\! s_{i_0...i_n} \!-\! \sum_{l=0}^n(-1)^l s_{ji_0...\hat{i_l}...i_n}|_{U_{i_0}\cap...\cap U_{i_n}} \\ 
		&= (\text{id}_{\check{\mathcal{C}}^n(\mathfrak{U},\mathcal{F})_x}-0)(s)_{i_0...i_n}\,, 
\end{align*}\end{small}
\endgroup
\!\!这里我们最后对指标 $l$ 作了平移并运用了反对称性。因此 $\text{id}_{\check{\mathcal{C}}^n(\mathfrak{U},\mathcal{F})_x}$ 同伦于零映射，故二者在上同调中相等，从而 $H^n(\check{\mathcal{C}}^\bullet(\mathfrak{U},\mathcal{F})_x)=H^n(\text{id}_{\check{\mathcal{C}}^\bullet(\mathfrak{U},\mathcal{F})_x})(H^n(\check{\mathcal{C}}^\bullet(\mathfrak{U},\mathcal{F})_x))\cong\{0\}$。最后，零调性蕴含 $\check{\mathcal{C}}^\bullet(\mathfrak{U},\mathcal{F})$ 的所述正合性。  
\end{proof}

% [Tutor Note: 这个证明是\'{e}tal\'{e} 上同调/超覆盖理论的雏形：$\check{\mathcal{C}}^\bullet(\mathfrak{U},\mathcal{F})$ 通过 $\varepsilon$ 给出的不是普通消解，而是"零调消解"，它不必内射。作者称其为内射消解，是因为下一段只需它是一个对 $\Gamma(X,\square)$ 零调的消解来比对层上同调；真正用到的性质是引理中的逐茎同伦，其中把指标 $j$ 插入的技术需要 $\mathfrak{U}$ 是开覆盖（而不是闭覆盖）。]
我们来到此前承诺的与层上同调的联系：

\begin{Pro}
设 $X$ 为带有开覆盖 $\mathfrak{U}=(U_i)_{i\in I}$ 的拓扑空间，$\mathcal{F}\in\textup{obj}(\mathsf{Sh}(X;\mathsf{Ab}))$。则对每个 $n\geq 0$ 存在自然映射
\begin{equation}\label{Cechsheafcohmap}
	h_n:\check{H}^n(\mathfrak{U},\mathcal{F})\rightarrow H^n(X,\mathcal{F})\,,
\end{equation}
且它对 $\mathcal{F}$ 具有函子性。
\end{Pro}

\begin{proof}
同调代数中一个标准结果（例如见 \cite[Theorem III.1.3]{[GM03]} 或 \cite[Theorem IV.4.4]{[HS97]}）断言：任意一对复形（一个零调、一个内射）的零次上同调群之间的任何态射，都可以唯一地（在同伦意义下）延拓为整个复形之间的上链映射。

在我们的情形，对给定的 $\mathcal{F}\in\textup{obj}(\mathsf{Sh}(X;\mathsf{Ab}))$，考虑 $\mathsf{Sh}(X;\mathsf{Ab})$ 中某个内射消解 $\mathcal{F}\rightarrow\mathcal{I}^\bullet$，以及引理 \ref{Cechinjres} 提供的零调内射消解 $\mathcal{F}\xrightarrow{\varepsilon}\check{\mathcal{C}}^\bullet(\mathfrak{U},\mathcal{F})$。由于引理 \ref{CechH0isglobalsect} 给出典范同构 $\check{H}^0(\mathfrak{U},\mathcal{F})\cong\Gamma(X,\mathcal{F})\cong H^0(X,\mathcal{F})\cong H^0(\Gamma(X,\mathcal{I}^\bullet))$（其中最后两个同构分别来自层上同调的定义与注 \ref{classicalderivedonresolutions}），由上述定理存在唯一（从而是自然的）上链映射 $\check{\mathcal{C}}^\bullet(\mathfrak{U},\mathcal{F})\rightarrow\mathcal{I}^\bullet$，它在 $\mathcal{F}$ 上还可取为恒等。于是依次对它应用函子 $\Gamma(X,\square):\mathsf{Kom}^+(\mathsf{Sh}(X;\mathsf{Ab}))\rightarrow \mathsf{Kom^+(Ab)}$ 与 $H^n:\mathsf{Kom^+(Ab)}\rightarrow\mathsf{Ab}$，我们得到 $\check{\mathcal{C}}^\bullet(\mathfrak{U},\mathcal{F})(X)=\check{C}^\bullet(\mathfrak{U},\mathcal{F})\rightarrow\Gamma(X,\mathcal{I}^\bullet)$，进而得到所求的 $h_n: \check{H}^n(\mathfrak{U},\mathcal{F})\rightarrow H^n(\Gamma(X,\mathcal{I}^\bullet))$ $\cong\mathsf{R}^n\Gamma(X,\square)(\mathcal{F})=H^n(X,\mathcal{F})$（再次使用 \eqref{classderinjres}）。此外，对 $\mathcal{F}$ 的函子性由构造可得。 
\end{proof}

\begin{Rem}
可以证明（借助适当的谱序列；见 \cite[Theorem III.8.3b]{[GM03]}）：若存在\textit{$\mathcal{F}$-零调的}开覆盖 $\mathfrak{U}=(U_i)_{i\in I}$——即对一切 $n>0$ 及所有可能的有限交 $V=U_{i_0}\cap...\cap U_{i_m}\subset X$（$m\geq 0$）都有 $H^n(V,\mathcal{F}|_V)\cong\{0\}$ 的覆盖——则上一命题中的典范映射 $h_n$ 实际上是同构！
\end{Rem}

不幸的是，人们一般只能与标准的、非零调的开覆盖打交道，从而失去了刚才强调的好处。不过，当我们处理诺特分离概形上的拟凝聚层时，情况会变得更有趣，这一点我们很快就会看到。我们先证明一个具有普遍价值的预备引理。

\begin{Lem}\label{Cechonflasqueistrivial}
设 $X$ 为带有开覆盖 $\mathfrak{U}=(U_i)_{i\in I}$ 的拓扑空间。设 $\mathcal{F}\in\textup{obj}(\mathsf{Sh}(X;\mathsf{Ab}))$ 为\textbf{松弛层}\index{sheaf!flasque}，即对每对开子集 $V\subset U\subset X$，其限制映射 $\rho_V^U:\mathcal{F}(U)\rightarrow\mathcal{F}(V)$ 都满射。则对一切 $n>0$ 有 $\check{H}^n(\mathfrak{U},\mathcal{F})\cong\{0\}$。
\end{Lem}

\begin{proof}(\textit{Sketch})
取引理 \ref{Cechinjres} 中的内射消解 $\mathcal{F}\rightarrow \check{\mathcal{C}}^\bullet(\mathfrak{U},\mathcal{F})$，并注意每个 $\check{\mathcal{C}}^n(\mathfrak{U},\mathcal{F})=\prod_{1\leq i_0<...<i_n\leq |I|}(\iota_{i_0...i_n})_*(\mathcal{F}|_{U_{i_0}\cap...\cap U_{i_n}})\in\text{obj}(\mathsf{Sh}(X;\mathsf{Ab}))$ 自身是松弛层：每个 $\mathcal{F}|_{U_{i_0}\cap...\cap U_{i_n}}$ 按定义松弛，其在连续映射（如 $\iota_{i_0...i_n}$）下的推前也松弛，最后松弛层的积也松弛。由于松弛层对 $\Gamma(X,\square)$ 零调（见 \cite[Remark III.2.5.1]{[Har77]}），我们的松弛消解 $\mathcal{F}\rightarrow \check{\mathcal{C}}^\bullet(\mathfrak{U},\mathcal{F})$（由 \eqref{classderinjres} 可用以计算层上同调）给出一切 $n>0$ 的平凡层上同调群 $H^n(X,\mathcal{F})$。因此对一切 $n>0$ 有 $\check{H}^n(\mathfrak{U},\mathcal{F})=H^n\big(\Gamma(X,\check{\mathcal{C}}^\bullet(\mathfrak{U},\mathcal{F}))\big)\cong H^n(X,\mathcal{F})\cong\{0\}$。 
\end{proof}

值得指出的是，任意环空间 $(X,\mathcal{S})$ 上的内射 $\mathcal{S}$-模都是松弛的，因此正如刚才结束的证明所利用的那样，对每个非零次数都诱导平凡的层上同调群。

\vspace*{0.3cm}
\noindent 最后，我们回到概形与拟凝聚/凝聚层。（在此我们温和地提醒读者注意注 \ref{disclaimer} 中的免责声明。）

\begin{Lem}\label{neq0sheafcohomonaffinesch}
设 $X=\textup{Spec}(R)$ 是环 $R$ 上的仿射概形（$R$ 取为诺特环，从而 $X$ 自身诺特），$\mathcal{O}_X$ 为其结构层。则对一切 $\mathcal{F}\in\textup{obj}(\mathsf{QCoh}(X;{}_{\mathcal{O}_X}\!\mathsf{Mod}))$ 与 $n>0$ 有 $H^n(X,\mathcal{F})\cong \{0\}$。（特别地，对仿射簇成立。）
\end{Lem}

\begin{proof}(\textit{Sketch})
固定某个 $\mathcal{F}\in\text{obj}(\mathsf{QCoh}(X;{}_{\mathcal{O}_X}\!\mathsf{Mod}))$，令 $M\coloneqq\Gamma(X,\mathcal{F})\in\text{obj}({}_R\mathsf{Mod})$。取 ${}_R\mathsf{Mod}$ 中的内射消解 $0\rightarrow M\rightarrow I^\bullet$（总存在），它在等价 \eqref{squaretildeequivalences} 下给出 $X$ 上拟凝聚层的正合序列 $0\rightarrow \widetilde{M}\cong\mathcal{F}\rightarrow\widetilde{I}^\bullet$。可以证明每个 $\widetilde{I}^\bullet$ 实际上是松弛的（\cite[Proposition III.3.4]{[Har77]}）。于是正如上一个证明所论证的，所得的松弛消解对每个 $n>0$ 给出平凡的层上同调群，而按定义 $H^0(X,\mathcal{F})=\Gamma(X,\mathcal{F})=M$。 
\end{proof}

\begin{Cor}\label{embedinflasque}
诺特概形 $(X,\mathcal{O}_X)$ 上的任意 $\mathcal{F}\in\textup{obj}(\mathsf{QCoh}(X;{}_{\mathcal{O}_X}\!\mathsf{Mod}))$ 都可嵌入一个松弛的 $\mathcal{G}\in\textup{obj}(\mathsf{QCoh}(X;{}_{\mathcal{O}_X}\!\mathsf{Mod}))$。
\end{Cor}

\begin{proof}
用有限多个仿射开集 $U_i=\text{Spec}(R_i)$ 覆盖 $X$，令 $\widetilde{M}_i\coloneqq\mathcal{F}|_{U_i}$ $\in\textup{obj}(\mathsf{QCoh}(U_i;{}_{\mathcal{O}_{U_i}}\!\mathsf{Mod}))$，其相应的 $M_i\in\text{obj}({}_{R_i}\mathsf{Mod})$ 对每个 $i$ 被假定嵌入内射模 $I_i\in\text{obj}({}_{R_i}\mathsf{Mod})$。记嵌入为 $\iota_i:U_i\hookrightarrow X$，令 $\mathcal{G}\coloneqq\bigoplus_{i\in I}(\iota_i)_*(\widetilde{I}_i)\in\text{obj}(\mathsf{QCoh}(X;{}_{\mathcal{O}_X}\!\mathsf{Mod}))$；它是松弛的，因为它是松弛层 $\widetilde{I}_i$ 的推前的直和（同样由 \cite[Proposition III.3.4]{[Har77]}），而它是拟凝聚的，因为它是 $\widetilde{I}_i\in\text{obj}(\mathsf{QCoh}(U_i;{}_{\mathcal{O}_{U_i}}\!\mathsf{Mod}))$ 的推前（推前保持拟凝聚性；见命题 \ref{pushpullcoherent}）的直和。

但按构造，对每个 $i$ 有 $\mathcal{F}|_{U_i}$ 单射入 $\widetilde{I}_i$，这又诱导出 $\mathcal{F}\hookrightarrow(\iota_i)_*(\widetilde{I}_i)$，最终给出单射 $\mathcal{F}\hookrightarrow\mathcal{G}$。故 $\mathcal{F}$ 可嵌入一个松弛的拟凝聚层。
\end{proof}

% [Tutor Note: 此处"嵌入内射模再推前求和"的构造是 Serre 型论证的标准套路：关键在于①诺特概形可用有限个仿射覆盖，②仿射概形上拟凝聚层与模的等价使内射模可转成松弛层，③松弛层在推前与直和下封闭。注意 $\mathcal{G}$ 只保证在 $\mathsf{QCoh}$ 中内射，一般不必在 $\mathsf{Coh}$ 中内射——这正是"先在 $\mathsf{QCoh}$ 中做导出函子、再限制回 $\mathsf{Coh}$"策略的根源。]
\begin{Rem}\label{QCohhasenoughinjectives}
我们有以下值得注意的推论：
\begin{itemize}[leftmargin=0.5cm]
	\item 引理 \ref{neq0sheafcohomonaffinesch} 可以加强为：诺特概形 $X$ 是仿射的，当且仅当对一切 $\mathcal{F}\in\text{obj}(\mathsf{QCoh}(X;{}_{\mathcal{O}_X}\!\mathsf{Mod}))$ 与 $n>0$ 有 $H^n(X,\mathcal{F})\cong \{0\}$。
	
	\item 上一个证明中的拟凝聚层 $\mathcal{G}$ 实际上是内射对象（在 $\mathsf{Sh}(X;{}_{\mathcal{O}_X}\!\mathsf{Mod})$ 中亦然），这意味着当 $X$ 是诺特概形时 $\mathsf{QCoh}(X;{}_{\mathcal{O}_X}\!\mathsf{Mod})$ 有足够多的内射对象（其实并不需要诺特性）！为给这份热情降温，我们指出它\textit{没有}足够多的投射对象，就像 $\mathsf{Sh}(X;{}_{\mathcal{O}_X}\!\mathsf{Mod})$ 也没有一样。
	
	\item 此外，在诺特概形 $(X,\mathcal{O}_X)$ 上，$\mathsf{QCoh}(X;{}_{\mathcal{O}_X}\!\mathsf{Mod})$ 的任意内射对象都是松弛的，这可用以证明层上同调也可以由 $\mathsf{R}\Gamma(X,\square):\mathsf{D^+}(\mathsf{QCoh}(X;{}_{\mathcal{O}_X}\!\mathsf{Mod}))\rightarrow \mathsf{D^+}(\mathsf{Vect}_\mathbb{K}^{\leq\infty})$ 来计算。  
\end{itemize}
\end{Rem}

\begin{Thm}\label{sheafcohomthroughCech}
设 $(X,\mathcal{O}_X)$ 为带有有序仿射开覆盖 $\mathfrak{U}=(U_i)_{i\in I}$ 的诺特分离概形，$\mathcal{F}\in\textup{obj}(\mathsf{QCoh}(X;{}_{\mathcal{O}_X}\!\mathsf{Mod}))$。则 \eqref{Cechsheafcohmap} 中的自然映射 $h_n$ 是群同构，也就是说，\v{C}ech 上同调计算层上同调：
\begin{equation}
	\check{H}^n(\mathfrak{U},\mathcal{F})\cong H^n(X,\mathcal{F})\,,
\end{equation}
对一切 $n\geq 0$ 成立。
\end{Thm}

\begin{proof}(\textit{Sketch})
对 $n=0$ 由引理 \ref{CechH0isglobalsect} 立得：$H^0(X,\mathcal{F})\cong\Gamma(X,\square)(\mathcal{F})\cong\check{H}^0(\mathfrak{U},\mathcal{F})$。对 $n>0$，首先把 $\mathcal{F}$ 嵌入松弛的 $\mathcal{G}\in\textup{obj}(\mathsf{QCoh}(X;{}_{\mathcal{O}_X}\!\mathsf{Mod}))$（由推论 \ref{embedinflasque} 可行），并通过该嵌入的余核——$\mathcal{O}_X$-模商层 $\mathcal{H}$——将其补成短正合序列 $0\rightarrow\mathcal{F}\rightarrow\mathcal{G}\rightarrow\mathcal{H}\rightarrow 0$。特别地，由命题 \ref{cokernelscoherent} 有 $\mathcal{H}\in\text{obj}(\mathsf{QCoh}(X;{}_{\mathcal{O}_X}\!\mathsf{Mod}))$。

现在注意到 $U_{i_0}\cap...\cap U_{i_n}\subset X$ 是仿射开子集（若 $X$ 不分离则不然！）。在其上，\cite[Proposition II.5.6]{[Har77]} 断言上述以拟凝聚层 $\mathcal{F}$ 开头的短正合序列诱导出短正合序列
\[
0\rightarrow\Gamma(U_{i_0}\cap...\cap U_{i_n},\mathcal{F})\rightarrow\Gamma(U_{i_0}\cap...\cap U_{i_n},\mathcal{G})\rightarrow\Gamma(U_{i_0}\cap...\cap U_{i_n},\mathcal{H})\rightarrow 0
\]
再对一切 $1\leq i_0<...< i_n\leq|I|$ 与任意 $n\geq 0$ 取积之后，就得到 $\mathsf{Kom^+(Ab)}$ 中的短正合序列 $0\rightarrow\check{C}^\bullet(\mathfrak{U},\mathcal{F})\rightarrow\check{C}^\bullet(\mathfrak{U},\mathcal{G})\rightarrow\check{C}^\bullet(\mathfrak{U},\mathcal{H})\rightarrow 0$。应用标准上同调函子 $H^n$ 便得到\footnote{警告：直接论证 $\mathcal{O}_X$-模的短正合序列 $0\rightarrow\mathcal{F}\rightarrow\mathcal{G}\rightarrow\mathcal{H}\rightarrow 0$ 在 \v{C}ech 上同调中诱导长正合序列是错误的，因为函子 $\check{H}^n(\mathfrak{U},\square)$ 一般不是上同调函子！（例如若 $\mathfrak{U}=\{X\}$，则 $\check{H}^0(\mathfrak{U},\square)=\Gamma(X,\square)$ 只是左正合！）} \v{C}ech 上同调群的长正合序列
\[
...\rightarrow\check{H}^n(\mathfrak{U},\mathcal{G})\cong\{0\}\rightarrow\check{H}^n(\mathfrak{U},\mathcal{H})\rightarrow \check{H}^{n+1}(\mathfrak{U},\mathcal{F})\rightarrow \check{H}^{n+1}(\mathfrak{U},\mathcal{G})\cong\{0\}\rightarrow...\,,
\]
其中我们用引理 \ref{Cechonflasqueistrivial} 推出 $\mathcal{G}$ 的各群平凡。于是对一切 $n>0$ 有 $\check{H}^n(\mathfrak{U},\mathcal{H})\cong \check{H}^{n+1}(\mathfrak{U},\mathcal{F})$。

另一方面，推论 \ref{longexactforsheafcohomology} 告诉我们 $0\rightarrow\mathcal{F}\rightarrow\mathcal{G}\rightarrow\mathcal{H}\rightarrow 0$ 直接给出层上同调群的长正合序列
\[
...\rightarrow H^n(X,\mathcal{G})\cong\{0\}\rightarrow H^n(X,\mathcal{H})\rightarrow H^{n+1}(X,\mathcal{F})\rightarrow H^{n+1}(X,\mathcal{G})\cong\{0\}\rightarrow...\,,
\]
其中我们利用了松弛层也使层上同调群平凡（由 \cite[Proposition III.2.5]{[Har77]}）。于是对一切 $n>0$ 有 $H^n(X,\mathcal{H})\cong H^{n+1}(X,\mathcal{F})$。

综上，由拟凝聚层 $\mathcal{H}$ 的 $n=0$ 情形得 $\check{H}^1(\mathfrak{U},\mathcal{F})\cong \check{H}^0(\mathfrak{U},\mathcal{H})\cong H^0(X,\mathcal{H})\cong H^1(X,\mathcal{F})$，命题遂由归纳法得证。
\end{proof}

\begin{Cor}\label{sheafcohomthroughCechRem}
设 $f:X\rightarrow Y$ 为诺特分离概形的仿射态射（即存在 $Y$ 的一个仿射开覆盖，其原像是 $X$ 的仿射开覆盖），$\mathcal{F}\in\textup{obj}(\mathsf{QCoh}(X;{}_{\mathcal{O}_X}\!\mathsf{Mod}))$。则对一切 $n\geq 0$ 有 $H^n(X,\mathcal{F})\cong H^n(Y,f_*\mathcal{F})$。
\end{Cor}

% [Tutor Note: 定理的证明有一个隐蔽但关键的技巧：先用松弛层 $\mathcal{G}$ 做成短正合序列，在 \v{C}ech 复形层面（$\mathsf{Kom^+}$ 中）得到短正合序列，此时取上同调才得到长正合序列；绝不能直接对 $\v{C}ech$ 上同调群用长正合序列（作者在脚注中特别警告了这一点）。另外"分离"假设在此不可省：交集 $U_{i_0}\cap...\cap U_{i_n}$ 仿射正依赖于分离性。]
最后，我们报告 Serre 的以下重要结果，它专门讨论射影概形的层上同调（见 \cite[Theorem III.5.2]{[Har77]}；我们在推论 \ref{globalsectforcoherent} 的证明中已经用过它），并推广了射影空间的情形（见 \cite[Theorem III.5.1]{[Har77]} 中的讨论）。

\begin{Thm}\label{sheafcohomforprojsch}
设 $(X,\mathcal{O}_X)$ 为诺特环 $R$ 上的射影概形，从而可逆扭层 $\mathcal{O}_X(1)$ 是非常丰富的（术语见命题 \ref{whenproper} 与注 \ref{asssheafofmodforproj}）。考虑任意 $\mathcal{F}\in\textup{obj}(\mathsf{Coh}(X;{}_{\mathcal{O}_X}\!\mathsf{Mod}))$。则对每个 $k\geq 0$，$H^k(X,\mathcal{F})$ 是有限生成的 $R$-模（特别地，$\Gamma(X,\mathcal{F})=\mathcal{F}(X)$ 也是）。

此外，存在某个 $n_0\equiv n_0(\mathcal{F})\in\mathbb{Z}$，使得对每个 $n\geq n_0$ 与 $k>0$ 有 $H^k(X,\mathcal{F}(n))\cong\{0\}$。
\end{Thm}

\subsection{凝聚层上的导出函子}\label{ch6.3}

从今往后，我们把精力集中于理解拟凝聚/凝聚层的导出范畴。我们一般在诺特概形上工作，因为在这样的概形上凝聚性是良态的。请注意，根据定理 \ref{imageunderXi}，以下分析对簇同样适用。接下来的材料取自 \cite[chapter 3]{[Huy06]}，\cite[chapter 3]{[Hoc19]} 中有更简短的综述。首先我们简化术语与记号。

\begin{Def}\label{dercatofsch}
设 $(X,\mathcal{O}_X)$ 为诺特概形，其相应的凝聚 $\mathcal{O}_X$-模范畴 $\mathsf{Coh}(X;{}_{\mathcal{O}_X}\!\mathsf{Mod})$ 由推论 \ref{QCohisabelian} 是阿贝尔范畴。则我们称其有界导出范畴
\[
\mathsf{D^b}(X)\coloneqq \mathsf{D^b}(\mathsf{Coh}(X;{}_{\mathcal{O}_X}\!\mathsf{Mod}))
\]
为（简言之）$X$ 的\textbf{导出范畴}\index{derived category!of a noetherian scheme}。当然，命题 \ref{derivedexists} 对 $\#=\emptyset,+,-$ 也构造了导出范畴 $\mathsf{D^\#\!}(X)\coloneqq\mathsf{D^\#\!}(\mathsf{Coh}(X;{}_{\mathcal{O}_X}\!\mathsf{Mod}))$。

我们称任意域 $\mathbb{K}$ 上的两个诺特概形 $X,Y\in\text{obj}(\mathsf{Sch}(\mathbb{K}))$ 是\textbf{导出等价的}\index{derived equivalent (noetherian schemes)}，若存在一个 $\mathbb{K}$-线性\footnote{当处理域 $\mathbb{K}$ 上的概形时，我们隐式地假定其导出范畴是 $\mathbb{K}$-线性的（其含义已在定义 \ref{addcat} 之后说明）。}的、正合的范畴等价 $\mathsf{D^b}(X)\rightarrow\mathsf{D^b}(Y)$（回想在导出范畴上，正合性总按定义 \ref{exactderivedfunc} 的意义理解）。
\end{Def}

不幸的是，$\mathsf{Coh}(X;{}_{\mathcal{O}_X}\!\mathsf{Mod})$ 既没有足够多的内射对象也没有足够多的投射对象，这使它在构造导出函子时显得不适用……但情形并非如此，因为我们可以先在更大的、由拟凝聚 $\mathcal{O}_X$-模构成的阿贝尔范畴 $\mathsf{QCoh}(X;{}_{\mathcal{O}_X}\!\mathsf{Mod})$ 中工作——它至少有足够多的内射对象（见注 \ref{QCohhasenoughinjectives}）以及合适的适配类——然后再限制到凝聚层！我们将反复利用这一点。

\begin{Rem}
正如注 \ref{QCohhasenoughinjectives} 所指出的，任意拟凝聚层都可嵌入一个内射 $\mathcal{O}_X$-模。此外，全阿贝尔子范畴 $\mathsf{QCoh}(X;{}_{\mathcal{O}_X}\!\mathsf{Mod})\subset\mathsf{Sh}(X;{}_{\mathcal{O}_X}\!\mathsf{Mod})$ 是厚的（见 \cite[Proposition II.5.7]{[Har77]}）。于是我们可以应用引理 \ref{abeliansubcatderivedequiv}，把 $\mathsf{D^\#\!}(\mathsf{QCoh}(X;{}_{\mathcal{O}_X}\!\mathsf{Mod}))$ 解释为 $\mathsf{D^\#\!}(\mathsf{Sh}(X;{}_{\mathcal{O}_X}\!\mathsf{Mod}))$ 中具有拟凝聚（标准）上同调的全子范畴。
\end{Rem}

% [Tutor Note: 记号 $\mathsf{D^b}(X)$ 专指凝聚层的有界导出范畴，与 $\mathsf{D^b}(\mathsf{QCoh})$、$\mathsf{D^b}(\mathsf{Sh})$ 都不同；后两者通过厚子范畴包含关系给出全子范畴（引理 \ref{boundedCohinQCoh} 说明 $\mathsf{D^b}(X)$ 恰是 $\mathsf{D^b}(\mathsf{QCoh})$ 中上同调凝聚的那些对象）。这一链式包含是本节全部构造的骨架：先在 $\mathsf{QCoh}$ 中造导出函子，再用上同调凝聚性把它压回 $\mathsf{D^b}(X)$。]
承上注，我们有（参看 \cite[Proposition 3.5]{[Huy06]}）：

\begin{Lem}\label{boundedCohinQCoh}
设 $X$ 为诺特概形。则自然包含函子
\[
\mathsf{I}_X:\mathsf{D^b}(X)\hookrightarrow\mathsf{D^b}(\mathsf{QCoh}(X;{}_{\mathcal{O}_X}\!\mathsf{Mod}))
\]
限制为 $\mathsf{D^b}(X)$ 与 $\mathsf{D^b}(\mathsf{QCoh}(X;{}_{\mathcal{O}_X}\!\mathsf{Mod}))$ 中具有凝聚 \textup(标准\textup)上同调的有界复形所构成的全\textup(三角\textup)子范畴之间的等价。
\end{Lem}

现在，我们逐一验证第 \ref{ch6.1} 节中所有导出函子在凝聚层层面的存在性，同时讨论真正的层复形在它们之下的像是什么。注意，对 $X$ 的假设会不时被细化（始终引用第 \ref{ch5.5} 节中、由图表 \eqref{bigdiagram} 总结的词汇），但最终这些假设总能被我们最终青睐的概形——即 Calabi--Yau 簇——所满足。

\begin{Pro}\label{globalsectforcoherentbdd}
设 $X\in\textup{obj}(\mathsf{Sch}(\mathbb{K}))$ 为射影的 \textup(或仅为固有的\textup)诺特概形。则 \eqref{globsectsfunconcoh} 中左正合整体截面函子 $\Gamma(X,\square):\mathsf{Coh}(X;{}_{\mathcal{O}_X}\!\mathsf{Mod})\rightarrow\mathsf{Vect}_\mathbb{K}$ 在有界复形上的\textup(正合\textup)右导出函子存在\footnote{注意这并非自动成立，因为如前所述 $\mathsf{Coh}(X;{}_{\mathcal{O}_X}\!\mathsf{Mod})$ 没有足够多的内射对象！}，并可表为正合函子的复合。
\end{Pro}

\begin{proof}(\textit{Sketch})
由推论 \ref{globalsectforcoherent}，并由于 $\mathsf{QCoh}(X;{}_{\mathcal{O}_X}\!\mathsf{Mod})$ 有足够多的内射对象，左正合的 $\Gamma(X,\square):$ $\mathsf{QCoh}(X;{}_{\mathcal{O}_X}\!\mathsf{Mod})\rightarrow\mathsf{Vect}_\mathbb{K}^{\leq\infty}$ 给出右导出函子 $\mathsf{R}\Gamma(X,\square):\mathsf{D^+}(\mathsf{QCoh}(X;{}_{\mathcal{O}_X}\!\mathsf{Mod}))\rightarrow\mathsf{D^+}(\mathsf{Vect}_\mathbb{K}^{\leq\infty})$。 

现在，Grothendieck 的一个定理（见 \cite[Theorem III.2.7]{[Har77]}）告诉我们：任意 $\mathcal{F}\in\text{obj}(\mathsf{QCoh}(X;{}_{\mathcal{O}_X}\!\mathsf{Mod}))$ 的层上同调群在 $n>\text{dim}(X)$ 时满足 $H^n(X,\mathcal{F})\cong\{0\}$。于是可应用命题 \ref{derivedfuncrestrictstobdd} 说明 $\mathsf{R}\Gamma(X,\square)$ 下降到有界复形的导出范畴，$\mathsf{R}\Gamma(X,\square):\mathsf{D^b}(\mathsf{QCoh}(X;{}_{\mathcal{O}_X}\!\mathsf{Mod}))\rightarrow\mathsf{D^b}(\mathsf{Vect}_\mathbb{K}^{\leq\infty})$。我们于是有图表
\begin{equation}\label{diagramforderivedfuncofbdd}
	\begin{tikzcd}
		\mathsf{D^+}(\mathsf{QCoh}(X;{}_{\mathcal{O}_X}\!\mathsf{Mod})) \arrow[rr, "{\mathsf{R}\Gamma(X,\square)}"] & & \mathsf{D^+}(\mathsf{Vect}_\mathbb{K}^{\leq\infty}) \\
		\mathsf{D^b}(\mathsf{QCoh}(X;{}_{\mathcal{O}_X}\!\mathsf{Mod})) \arrow[rr, "{\mathsf{R}\Gamma(X,\square)}"']\arrow[u, hook] & & \mathsf{D^b}(\mathsf{Vect}_\mathbb{K}^{\leq\infty})\arrow[u, hook] \\
		\mathsf{D^b}(X) \arrow[u, hook, "\mathsf{I}_X"]\arrow[rr, dashed] & & \mathsf{D^b}(\mathsf{Vect}_\mathbb{K})\arrow[u, hook]
	\end{tikzcd}\quad,
\end{equation}
其中 $\mathsf{I}_X$ 仍表示自然包含函子。但由引理 \ref{boundedCohinQCoh}，后者实际上落在具有凝聚上同调的有界复形构成的全子范畴中，而该子范畴在中间那个 $\mathsf{R}\Gamma(X,\square)$ 下又映到有限维 $\mathbb{K}$-向量空间的有界复形，于是我们看到实际上可以绕图表的下方形走，得到所画的复合虚线箭头。又因为任意凝聚模的其余层上同调群都是有限维的（由定理 \ref{sheafcohomforprojsch}），且 $\mathsf{Vect}_\mathbb{K}\subset \mathsf{Vect}_\mathbb{K}^{\leq\infty}$ 显然是厚的，对这条虚线箭头再用一次命题 \ref{derivedfuncrestrictstobdd} 就把它变成所求的函子 $\mathsf{D^b}(X)\rightarrow \mathsf{D^b}(\mathsf{Vect}_\mathbb{K})$。   	     
\end{proof}

% [Tutor Note: 作者对右导出函子用了"正合"（exact）一词，这是导出范畴语境下的行话：$\mathsf{R}\Gamma$ 作为三角范畴之间的函子是正合的，但它当然不是阿贝尔范畴层面的正合函子。整个证明是"先在大范畴 $\mathsf{QCoh}$ 里造 $\mathsf{R}\Gamma$，再用 Grothendieck 消没定理把复形压到有界、最后用厚子范畴把像压回有限维向量空间"的三步走。]
于是对任意 $\mathcal{F}^\bullet\in\text{obj}(\mathsf{D^b}(X))$，我们可以定义 $H^n(X,\mathcal{F}^\bullet)\coloneqq\mathsf{R}^n\Gamma(X,\mathcal{F}^\bullet)=$ $H^n(\mathsf{R}\Gamma(X,\mathcal{F}^\bullet))\in\text{obj}(\mathsf{Vect}_\mathbb{K})$，从而由向量空间的自然分裂得到：
\[
\mathsf{R}\Gamma(X,\mathcal{F}^\bullet)\cong\bigoplus_{n\in\mathbb{Z}}\mathsf{R}^n\Gamma(X,\mathcal{F}^\bullet)[-n]=\bigoplus_{n\in\mathbb{Z}}H^n(X,\mathcal{F}^\bullet)[-n]\in\text{obj}(\mathsf{D^+}(\mathsf{Vect}_\mathbb{K}))\,.
\]

\vspace*{0.3cm}
\noindent 让我们转向更一般的推前函子的情形。

\begin{Pro}
设 $f:X\rightarrow Y$ 为诺特概形的射影 \textup(或仅为固有的\textup)态射。则命题 \ref{pushpullcoherent} 中左正合推前函子 $f_*:\mathsf{Coh}(X;{}_{\mathcal{O}_X}\!\mathsf{Mod})\rightarrow\mathsf{Coh}(Y;{}_{\mathcal{O}_Y}\!\mathsf{Mod})$ 在有界复形上的右导出函子存在，并可表为正合函子的复合。\textup(取 $Y=\textup{Spec}(\mathbb{K})$ 且 $f$ 为 $X\in\textup{obj}(\mathsf{Sch}(\mathbb{K}))$ 的结构态射，即得命题 \ref{globalsectforcoherentbdd}。\textup)
\end{Pro}

\begin{proof}(\textit{Sketch})
我们只需模仿整体截面函子的情形之证明。于是由命题 \ref{pushpullcoherent}，在拟凝聚层上可以构造右导出函子 $\mathsf{R}f_*:\mathsf{D^+}(\mathsf{QCoh}(X;{}_{\mathcal{O}_X}\!\mathsf{Mod}))\rightarrow\mathsf{D^+}(\mathsf{QCoh}(Y;{}_{\mathcal{O}_Y}\!\mathsf{Mod}))$。

现在，Grothendieck 的定理自然地推广为：对任意 $\mathcal{F}\in\text{obj}(\mathsf{QCoh}(X;{}_{\mathcal{O}_X}\!\mathsf{Mod}))$，若 $n>\text{dim}(X)$ 则 $\mathsf{R}^nf_*(\mathcal{F})\cong\{0\}$（$f$ 只需是诺特概形的一般态射即可；这一结果例如可由 \cite[Proposition III.8.1]{[Har77]} 推出）。对命题 \ref{derivedfuncrestrictstobdd} 的第一次调用随之把 $\mathsf{R}f_*$ 特化为 $\mathsf{R}f_*:\mathsf{D^b}(\mathsf{QCoh}(X;{}_{\mathcal{O}_X}\!\mathsf{Mod}))\rightarrow\mathsf{D^b}(\mathsf{QCoh}(Y;{}_{\mathcal{O}_Y}\!\mathsf{Mod}))$。

但若我们处理凝聚层 $\mathcal{F}\in\text{obj}(\mathsf{Coh}(X;{}_{\mathcal{O}_X}\!\mathsf{Mod}))$ 且 $f$ 如假设，则 $\mathsf{R}^nf_*(\mathcal{F})$ 对一切 $n\geq 0$ 也是凝聚的（见 \cite[Theorem III.8.8]{[Har77]}）。结合上述，我们确实可以通过复合构造 $\mathsf{R}f_*:\mathsf{D^b}(X)\rightarrow\mathsf{D^b}(Y)$，从而得到图表 \eqref{diagramforderivedfuncofbdd} 的推广版本中的虚线箭头（再次使用引理 \ref{boundedCohinQCoh} 与命题 \ref{derivedfuncrestrictstobdd}）。
\end{proof}

于是对任意 $\mathcal{F}^\bullet\in\text{obj}(\mathsf{D^b}(X))$，我们可以定义高阶直接像
\begin{equation}
\mathsf{R}^nf_*(\mathcal{F}^\bullet)\coloneqq H^n(\mathsf{R}f_*(\mathcal{F}^\bullet))\in\text{obj}(\mathsf{Coh}(Y;{}_{\mathcal{O}_Y}\!\mathsf{Mod}))
\end{equation}
\big(当 $Y=\text{Spec}(\mathbb{K})$ 且 $f$ 是 $X\in\text{obj}(\mathsf{Sch}(\mathbb{K}))$ 的结构态射时，它当然等于 $H^n(X,\mathcal{F}^\bullet)$\big)。

\begin{Rem}
推前运算关于复合表现良好：对诺特概形的相容态射 $f:X\rightarrow Y$ 与 $g:Y\rightarrow Z$，有 $(g\circ f)_*=g_*\circ f_*:\mathsf{QCoh}(X;{}_{\mathcal{O}_X}\!\mathsf{Mod})\rightarrow\mathsf{QCoh}(Z;{}_{\mathcal{O}_Z}\!\mathsf{Mod})$。于是可以应用命题 \ref{derivedfunccompo} 证明
\[
\mathsf{R}(g\circ f)_*\cong\mathsf{R}g_*\circ\mathsf{R}f_*:\mathsf{D^+}(\mathsf{QCoh}(X;{}_{\mathcal{O}_X}\!\mathsf{Mod}))\rightarrow\mathsf{D^+}(\mathsf{QCoh}(Z;{}_{\mathcal{O}_Z}\!\mathsf{Mod}))\,.
\]
（完全像命题 \ref{sheafcohomologyisabelian} 的证明那样，取内射拟凝聚层、相应地取松弛拟凝聚层作为适配类。）当 $Z=\text{Spec}(\mathbb{K})$ 且 $g$ 是 $Y$ 的结构态射时它最为有用：
\begin{equation}\label{globalsectcompo}
	\mathsf{R}\Gamma(X,\square)\cong\mathsf{R}\Gamma(Y,\square)\circ\mathsf{R}f_*:\mathsf{D^+}(\mathsf{QCoh}(X;{}_{\mathcal{O}_X}\!\mathsf{Mod}))\rightarrow\mathsf{D^+}(\mathsf{Vect}_\mathbb{K}^{\leq\infty})\,.
\end{equation}
\end{Rem}

\vspace*{0.3cm}
\noindent 根据注 \ref{desirablecoherentprop}，局部 $\mathcal{H}om$ 是另一类有趣函子的来源。首先注意，我们可以把 \eqref{locHomonqcoh} 推广为左正合函子
\begin{equation}\label{locHomoncomplexes}
\begin{aligned}
	&\mathcal{H}om_{\mathcal{O}_X}^\bullet(\mathcal{F}^\bullet,\square):\mathsf{H^+}(\mathsf{QCoh}(X;{}_{\mathcal{O}_X}\!\mathsf{Mod}))\rightarrow\mathsf{H^+}(\mathsf{QCoh}(X;{}_{\mathcal{O}_X}\!\mathsf{Mod}))\,, \\
	&\mathcal{H}om_{\mathcal{O}_X}^n(\mathcal{F}^\bullet,\mathcal{G}^\bullet)\coloneqq \prod_{k\in\mathbb{Z}}\mathcal{H}om_{\mathcal{O}_X}(\mathcal{F}^k,\mathcal{G}^{n+k})\in\text{obj}(\mathsf{QCoh}(X;{}_{\mathcal{O}_X}\!\mathsf{Mod})) \\
	&\text{and}\quad d^n\coloneqq d_\mathcal{G}^n -(-1)^n d_\mathcal{F}^n:\mathcal{H}om_{\mathcal{O}_X}^n(\mathcal{F}^\bullet,\mathcal{G}^\bullet)\rightarrow \mathcal{H}om_{\mathcal{O}_X}^{n+1}(\mathcal{F}^\bullet,\mathcal{G}^\bullet)\,,
\end{aligned}
\end{equation}
其中 $\mathcal{F}^\bullet$ 是 $\mathsf{H^-}(\mathsf{QCoh}(X;{}_{\mathcal{O}_X}\!\mathsf{Mod}))$ 中任意真正的拟凝聚层复形。 

\begin{Lem}
设 $X$ 为诺特概形且是正则的，按定义 \ref{smoothsch} 的意义 \textup{\big(}这适用于任意光滑的 $X\in\textup{obj}(\mathsf{Sch}(\mathbb{K}))$\textup{\big)}。则对任意 $\mathcal{F}\in\textup{obj}(\mathsf{Coh}(X;{}_{\mathcal{O}_X}\!\mathsf{Mod}))$，右导出函子 $\mathsf{R}\mathcal{H}om_{\mathcal{O}_X}(\mathcal{F},\square):\mathsf{D^b}(X)\rightarrow\mathsf{D^b}(X)$ 存在，并推广为正合函子
\[
\mathsf{R}\mathcal{H}om_{\mathcal{O}_X}^\bullet(\square,\square):\mathsf{D^b}(X)^\textup{opp}\times\mathsf{D^b}(X)\rightarrow\mathsf{D^b}(X)\,.
\] 
\end{Lem}

\begin{proof}(\textit{Sketch})
事实上，对诺特概形 $X$ 上的任意 $\mathcal{F}\in\text{obj}(\mathsf{QCoh}(X;{}_{\mathcal{O}_X}\!\mathsf{Mod}))$，我们得到右导出函子
\[
\mathsf{R}\mathcal{H}om_{\mathcal{O}_X}(\mathcal{F},\square):\mathsf{D^+}(\mathsf{QCoh}(X;{}_{\mathcal{O}_X}\!\mathsf{Mod}))\rightarrow\mathsf{D^+}(\mathsf{QCoh}(X;{}_{\mathcal{O}_X}\!\mathsf{Mod}))\,.
\]
若进一步假定 $\mathcal{F}$ 凝聚，则后者可以限制为右导出函子 $\mathsf{R}\mathcal{H}om_{\mathcal{O}_X}(\mathcal{F},\square):\mathsf{D^+}(X)\rightarrow\mathsf{D^+}(X)$；但要得到正合函子 $\mathsf{D^b}(X)\rightarrow\mathsf{D^b}(X)$，我们实际上需要 $X$ 是正则概形。

另一方面，从 $\mathsf{H^-}(\mathsf{QCoh}(X;{}_{\mathcal{O}_X}\!\mathsf{Mod}))$ 中的某个 $\mathcal{F}^\bullet$ 出发，可以验证内射层复形构成的类是 $\mathcal{H}om_{\mathcal{O}_X}^\bullet(\mathcal{F}^\bullet,\square)$ 的适配类，从而构造出 $\mathsf{R}\mathcal{H}om_{\mathcal{O}_X}^\bullet(\mathcal{F}^\bullet,\square):\mathsf{D^+}(\mathsf{QCoh}(X;{}_{\mathcal{O}_X}\!\mathsf{Mod}))\rightarrow\mathsf{D^+}(\mathsf{QCoh}(X;{}_{\mathcal{O}_X}\!\mathsf{Mod}))$，由此得到一个左变元在 $\mathsf{D^-}(\mathsf{QCoh}(X;{}_{\mathcal{O}_X}\!\mathsf{Mod}))^\text{opp}$ 中的双函子。最后，若 $X$ 正则，我们就得到凝聚层有界复形上所求的正合函子。
更多细节请读者参阅 \cite{[Huy06]}。
\end{proof}

对 $\mathcal{F},\mathcal{G}\in\text{obj}(\mathsf{QCoh}(X;{}_{\mathcal{O}_X}\!\mathsf{Mod}))$，此时（与定义 \ref{localExt} 一致地）如下定义局部 $\mathcal{E}xt$-层是有意义的：
\begin{equation}
\mathcal{E}xt_{\mathcal{O}_X}^n(\mathcal{F},\mathcal{G})\coloneqq\mathsf{R}^n\mathcal{H}om_{\mathcal{O}_X}(\mathcal{F},\mathcal{G})\in\text{obj}(\mathsf{QCoh}(X;{}_{\mathcal{O}_X}\!\mathsf{Mod}))\,,
\end{equation}
并且只要 $\mathcal{F},\mathcal{G}$ 都凝聚，它也是凝聚的。类似地，设 $\mathcal{F}^\bullet$ 属于 $\mathsf{D^-}(\mathsf{QCoh}(X;{}_{\mathcal{O}_X}\!\mathsf{Mod}))$、$\mathcal{G}^\bullet$ 属于 $\mathsf{D^+}(\mathsf{QCoh}(X;{}_{\mathcal{O}_X}\!\mathsf{Mod}))$，我们令
\begin{equation}
\mathcal{E}xt_{\mathcal{O}_X}^n(\mathcal{F}^\bullet,\mathcal{G}^\bullet)\coloneqq\mathsf{R}^n\mathcal{H}om_{\mathcal{O}_X}^\bullet(\mathcal{F}^\bullet,\mathcal{G}^\bullet)\in\text{obj}(\mathsf{QCoh}(X;{}_{\mathcal{O}_X}\!\mathsf{Mod}))\,,
\end{equation}
并且只要 $\mathcal{F}^\bullet,\mathcal{G}^\bullet\in\text{obj}(\mathsf{D^b}(X))$ 它就是凝聚的。

\begin{Rem}\label{HomExtRem}
以下是算子 $\mathcal{H}om$、$\mathcal{E}xt$、Hom 与 Ext 之间的一些值得注意的关系。
\begin{itemize}[leftmargin=0.5cm]
	\item 在附加假设 $X\in\text{obj}(\mathsf{Sch}(\mathbb{K}))$ 射影之下，仿照上述论证，由 $\mathsf{H^+}(\mathsf{QCoh}(X;{}_{\mathcal{O}_X}\!\mathsf{Mod}))$ 上按 \eqref{locHomoncomplexes} 完全相同地定义的内 Hom 左正合函子 $\text{Hom}_{\mathcal{O}_X}^\bullet(\mathcal{F}^\bullet,\square)$，我们可以构造正合双函子 $\mathsf{R}\text{Hom}_{\mathcal{O}_X}^\bullet\!(\square,\square)\!:\!\mathsf{D^b}(X)^\textup{opp}\times\mathsf{D^b}(X)\rightarrow\mathsf{D^b}(\mathsf{Vect}_\mathbb{K})$，它推广了我们从 \eqref{locHomoncoh} 会得到的右导出函子 $\mathsf{R}\text{Hom}_{\mathcal{O}_X}\!(\mathcal{F},\square)$。于是对任意 $\mathcal{F}^\bullet,\mathcal{G}^\bullet\in\text{obj}(\mathsf{D^b}(X))$，我们可以令 
	\begin{small}\begin{equation}\label{ExtRHomforcomplexes}
			\mkern-12mu\text{Ext}_{\mathcal{O}_X}^n(\mathcal{F}^\bullet,\mathcal{G}^\bullet)\!\coloneqq\!\mathsf{R}^n\text{Hom}_{\mathcal{O}_X}^\bullet(\mathcal{F}^\bullet,\mathcal{G}^\bullet)\!\cong\!\text{Hom}_{\mathsf{D^b}(X)}(\mathcal{F}^\bullet,\mathcal{G}^\bullet[n])\!\in\!\text{obj}(\mathsf{Vect}_\mathbb{K})\,,
	\end{equation}\end{small}
	\!\!这完全符合我们对 Ext-模的标准定义 \eqref{Extmodules}。注意由关于 $X$ 的假设与定理 \ref{sheafcohomforprojsch}，$\text{Ext}_{\mathcal{O}_X}^n\!(\mathcal{F}^\bullet,\mathcal{G}^\bullet)$ 是有限维的。\footnote{因此这对射影簇 $X$ 也成立；若假定它光滑，则可以证明当 $n>\text{dim}(X)$ 且 $\mathcal{F},\mathcal{G}$ 凝聚时 $\text{Ext}_{\mathcal{O}_X}^n\!(\mathcal{F},\mathcal{G})\cong\{0\}$（从而对同调维数 \eqref{homologicaldimension} 有 $\text{dim}_h(\mathsf{Coh}(X;{}_{\mathcal{O}_X}\!\mathsf{Mod}))=\text{dim}(X)$），并且对凝聚层的有界复形 $\mathcal{F}^\bullet,\mathcal{G}^\bullet$ 有 $\mathsf{R}\text{Hom}_{\mathcal{O}_X}^\bullet(\mathcal{F}^\bullet,\mathcal{G}^\bullet)\in\text{obj}(\mathsf{D^b(Ab)})$。} 此外，对任意 $\mathcal{F}^\bullet\in\textup{obj}(\mathsf{D^-}(X))$ 有
	\begin{equation}\label{locHomcompo}
		\mathsf{R}\Gamma(X,\square)\circ\mathsf{R}\mathcal{H}om_{\mathcal{O}_X}\!(\mathcal{F}^\bullet,\square)\cong\mathsf{R}\textup{Hom}_{\mathcal{O}_X}\!(\mathcal{F}^\bullet,\square)\!:\mathsf{D^+}(X)\!\rightarrow~{\!\mathsf{D^+}(X)}\,.
	\end{equation} 
	
	\item 设 $X$ 为诺特概形，$\mathcal{F}\in\text{obj}(\mathsf{Coh}(X;{}_{\mathcal{O}_X}\!\mathsf{Mod}))$，$\mathcal{G}\in\text{obj}(\mathsf{Sh}(X;{}_{\mathcal{O}_X}\mathsf{Mod}))$。此时局部 $\mathcal{E}xt$-层与标准 $\text{Ext}$-模在茎的层面上有如下关系
	\[
	\mathcal{E}xt_{\mathcal{O}_X}^n(\mathcal{F},\mathcal{G})_x\cong\text{Ext}_{\mathcal{O}_{X,x}}^n(\mathcal{F}_x,\mathcal{G}_x)\in\text{obj}({}_{\mathcal{O}_{X,x}}\!\mathsf{Mod})\,,
	\]
	对一切 $x\in X$ 与 $n\geq 0$ 成立（见 \cite[Proposition III.6.8]{[Har77]}）。
	\newline 进一步，设 $X=\text{Proj}(S)$ 是某个诺特分次环 $S$ 上的射影概形，$\mathcal{F},\mathcal{G}\in$ $\text{obj}(\mathsf{Coh}(X;{}_{\mathcal{O}_X}\!\mathsf{Mod}))$，则可以证明 Ext-模"最终"成为 $\mathcal{E}xt$-模的整体截面模（参看 \cite[Proposition III.6.9]{[Har77]}）。
	
	\item 给定 $\mathsf{D^-}(\mathsf{QCoh}(X;{}_{\mathcal{O}_X}\!\mathsf{Mod}))$ 中任意 $\mathcal{F}^\bullet$，我们定义其对偶复形为 
	\begin{equation}
		(\mathcal{F}^\bullet)^\vee\coloneqq\mathsf{R}\mathcal{H}om(\mathcal{F}^\bullet,\mathcal{O}_X)\in\text{obj}\big(\mathsf{D^+}(\mathsf{QCoh}(X;{}_{\mathcal{O}_X}\!\mathsf{Mod}))\big)
	\end{equation}
	（假定 $\mathcal{O}_X$ 集中在次数 0）。若 $X$ 是正则概形，则对任意 $\mathcal{F}^\bullet\in\text{obj}(\mathsf{D^b}(X))$，$(\mathcal{F}^\bullet)^\vee$ 实际上属于 $\mathsf{D^b}(X)$。
\end{itemize}
\end{Rem}

% [Tutor Note: 本注第一个等式组的关键在于第二个同构：它把导出范畴中的 Hom 与 $\text{Ext}$-模直接等同，这是同调镜像对称中「B-侧态射由 Ext 群给出」的技术根据。另需留意本注末尾那个等式的右端写作 $\mathsf{D^+}(X)$，按内容应是取值于向量空间的导出范畴；作者此处记号有笔误。]
\vspace*{0.7cm}
\noindent 我们把注意力转向拟凝聚/凝聚层上的右正合函子，从张量积开始。正如局部 $\mathcal{H}om$ 的情形，我们可以把 \eqref{tensoroncoh} 中的 $\mathcal{F}\otimes_{\mathcal{O}_X}\!\square$ 推广为右正合函子
\begin{equation}
\begin{aligned}
	&\mkern-18mu\mathcal{F}^\bullet\otimes_{\mathcal{O}_X}\!\square:\mathsf{H^-}(\mathsf{Coh}(X;{}_{\mathcal{O}_X}\!\mathsf{Mod}))\rightarrow\mathsf{H^-}(\mathsf{Coh}(X;{}_{\mathcal{O}_X}\!\mathsf{Mod}))\,, \\
	&\mkern-18mu(\mathcal{F}^\bullet\!\otimes_{\mathcal{O}_X}\!\mathcal{G}^\bullet)^n\!\coloneqq \mkern-12mu\bigoplus_{k+l=n\in\mathbb{Z}}\mkern-12mu\mathcal{F}^k\!\otimes_{\mathcal{O}_X}\!\mathcal{G}^l \;\;\text{and}\;\; d^n\!\coloneqq\!d_\mathcal{F}^n\otimes\text{id}_\mathcal{G}^n \!+\!(-1)^n \text{id}_\mathcal{F}^n\otimes d_\mathcal{G}^n\,,
\end{aligned}
\end{equation}
其中 $\mathcal{F}^\bullet$ 为 $\mathsf{H^-}(\mathsf{Coh}(X;{}_{\mathcal{O}_X}\!\mathsf{Mod}))$ 中任意对象；对右正合的 $\square\!\otimes_{\mathcal{O}_X}\mathcal{G}^\bullet$ 类似。

\begin{Pro}\label{tensorbifunconbdd}
设 $X\in\textup{obj}(\mathsf{Sch}(\mathbb{K}))$ 光滑且射影。给定任意 $\mathcal{F}\in\textup{obj}(\mathsf{Coh}(X;\mathsf{Mod}_{\mathcal{O}_X}))$ 与 $\mathcal{G}\in\textup{obj}(\mathsf{Coh}(X;{}_{\mathcal{O}_X}\!\mathsf{Mod}))$，左导出函子 $\mathcal{F}\overset{\mathsf{L}}{\otimes}_{\mathcal{O}_X}\!\square\coloneqq\mathsf{L}(\mathcal{F}\otimes_{\mathcal{O}_X}\!\square):\mathsf{D^b}(X)\rightarrow\mathsf{D^b}(X)$ 与 $\square\overset{\mathsf{L}}{\otimes}_{\mathcal{O}_X}\mathcal{G}:$ $\mathsf{D^b}(X)\rightarrow\mathsf{D^b}(X)$ 存在，并推广为正合函子
\[
\square\overset{\mathsf{L}}{\otimes}_{\mathcal{O}_X}\!\square:\mathsf{D^b}(X)\times\mathsf{D^b}(X)\rightarrow\mathsf{D^b}(X)
\]
\textup(它在两个变元上既交换又结合，如同标准的张量积\textup)。
\end{Pro}

\begin{proof}(\textit{Sketch})
设 $\mathcal{F}\in\textup{obj}(\mathsf{Coh}(X;\mathsf{Mod}_{\mathcal{O}_X}))$。注 \ref{localExtRem} 的第一条提示我们：任意 $\mathcal{E}\in\textup{obj}(\mathsf{Coh}(X;{}_{\mathcal{O}_X}\!\mathsf{Mod}))$ 都容许一个由局部自由层构成的投射消解 $\mathcal{L}^\bullet\rightarrow\mathcal{E}$：事实上，诺特环上射影概形上的任意凝聚层都可写成有限秩局部自由层的商（见 \cite[Corollary II.5.18]{[Har77]}）。这就切合于 $\mathsf{Coh}(X;{}_{\mathcal{O}_X}\!\mathsf{Mod})$ 中局部自由层构成的类适配于 $\mathcal{F}\otimes_{\mathcal{O}_X}\!\square$，从而证明我们可以构造左导出函子 $\mathsf{L}(\mathcal{F}\otimes_{\mathcal{O}_X}\!\square):\mathsf{D^-}(X)\rightarrow\mathsf{D^-}(X)$。

但由于 $X$ 是某个有限维 $\text{dim}(X)$ 的光滑概形，可以验证 $\mathcal{E}$ 的上述任一消解长度也是有限的；用它来计算经典导出函子（这一步由 \eqref{classderprojres} 是合法的）便得到对一切 $-n<-\text{dim}(X)$ 有 $\text{Tor}_n^{\mathcal{O}_X}(\mathcal{F},\mathcal{G})=\mathsf{H}^{-n}(\mathcal{F}\overset{\mathsf{L}}{\otimes}_{\mathcal{O}_X}\mathcal{G})\cong\{0\}$（回想例 \ref{derivedtensorfunc} 中的记号）。于是可以应用通常的命题 \ref{derivedfuncrestrictstobdd} 把张量函子限制为所求的正合函子 $\mathcal{F}\,\overset{\mathsf{L}}{\otimes}_{\mathcal{O}_X}\!\square:\mathsf{D^b}(X)\rightarrow\mathsf{D^b}(X)$。对 $\square\;\overset{\mathsf{L}}{\otimes}_{\mathcal{O}_X}\!\mathcal{G}$ 的论证是类似的。

更一般地，若 $\mathcal{F}^\bullet\in\text{obj}\big(\mathsf{H^-}(\mathsf{Coh}(X;\mathsf{Mod}_{\mathcal{O}_X}))\big)$，则同样可证局部自由层的复形构成的类适配于 $\mathcal{F}^\bullet\otimes_{\mathcal{O}_X}\!\square$，从而得到其左导出函子 $\mathcal{F}^\bullet\overset{\mathsf{L}}{\otimes}_{\mathcal{O}_X}\!\square:\mathsf{D^-}(X)\rightarrow\mathsf{D^-}(X)$；对右变元亦同。由于 $X$ 光滑，所得双函子再次下降到有界导出范畴 $\mathsf{D^b}(X)$。          
\end{proof}

与例 \ref{derivedtensorfunc} 一致，我们令
\begin{equation}
\text{Tor}_n^{\mathcal{O}_X}(\mathcal{F}^\bullet,\mathcal{G}^\bullet)\coloneqq H^{-n}(\mathcal{F}^\bullet\overset{\mathsf{L}}{\otimes}_{\mathcal{O}_X}\mathcal{G}^\bullet)\in\text{obj}(\mathsf{Coh}(X;{}_{\mathcal{O}_X}\!\mathsf{Mod}))\,,
\end{equation}
对任意 $\mathcal{F}^\bullet,\mathcal{G}^\bullet\in\text{obj}(\mathsf{D^b}(X))$ 成立。

\vspace*{0.3cm}
\noindent 我们这班人马中剩下的成员是拉回函子；由命题 \ref{pushpullcoherent}，我们知道它的要求比推前略低一些。

\begin{Cor}
设 $f:X\rightarrow Y$ 为光滑诺特概形的态射。则命题 \ref{pushpullcoherent} 中右正合拉回函子 $f^*:\mathsf{Coh}(Y;{}_{\mathcal{O}_Y}\!\mathsf{Mod})\rightarrow\mathsf{Coh}(X;{}_{\mathcal{O}_X}\!\mathsf{Mod})$ 在有界复形上的左导出函子 $\mathsf{L}f^*:\mathsf{D^b}(Y)\rightarrow\mathsf{D^b}(X)$ 存在，并可表为正合函子的复合。
\end{Cor}

\begin{proof}
这是 $f^*$ 的定义与命题 \ref{tensorbifunconbdd} 的直接推论；后者可以证明在这一更一般的设定下也成立（要求 $\mathbb{K}$ 上的光滑射影概形只是为了简便）。事实上，正合函子
\[
\mathsf{L}f^*\coloneqq \mathcal{O}_X\overset{\mathsf{L}}{\otimes}_{f^{-1}\mathcal{O}_Y}\!f^{-1}(\square):\mathsf{D^-}(Y)\rightarrow\mathsf{D^-}(X)
\]
下降到相应的有界复形导出范畴。
\end{proof} 

对任意 $\mathcal{G}^\bullet\in\text{obj}(\mathsf{D^b}(Y))$，我们令
\begin{equation}
\mathsf{L}^nf^*(\mathcal{F}^\bullet)\coloneqq H^n(\mathsf{L}f^*(\mathcal{F}^\bullet))\in\text{obj}(\mathsf{Coh}(X;{}_{\mathcal{O}_X}\!\mathsf{Mod}))\,.
\end{equation}

最后，我们列出一份颇为慷慨的恒等式清单，它们在操作凝聚层的有界复形时很有用。其证明留给一路活到这里的读者。建议：只需把分析化归到局部自由层——它们构成凝聚层上函子的适配类——从而回忆注 \ref{locallyfreeRem} 以及关于 $\mathsf{Sh}^\mathsf{lf}_n(X;{}_{\mathcal{O}_X}\!\mathsf{Mod})$ 的其他等式即可。

\begin{Pro}\label{lotofequalities}
以下始终设 $X,Y$ 为诺特概形，$f:X\rightarrow Y$ 为概形态射，$\mathcal{F}^\bullet, \mathcal{F}_i^\bullet\in\textup{obj}(\mathsf{D^b}(X))$，$\mathcal{G}^\bullet, \mathcal{G}_i^\bullet\in\textup{obj}(\mathsf{D^b}(Y))$ \textup(进一步的假设逐条在括号中说明\textup)。则以下关系成立：\begin{small}
	\begingroup
	\allowdisplaybreaks
	\begin{align}
		&a)\mkern+6mu\mathsf{R}f_*(\mathcal{F}^\bullet)\overset{\mathsf{L}}{\otimes}_{\mathcal{O}_Y}\!\mathcal{G}^\bullet\cong\mathsf{R}f_*(\mathcal{F}^\bullet\overset{\mathsf{L}}{\otimes}_{\mathcal{O}_X}\!\mathsf{L}f^*(\mathcal{G}^\bullet)) &\mkern-50mu\big(f\text{ proper},\, X,Y\!\in\!\textup{obj}(\mathsf{Sch}(\mathbb{K}))\big) \nonumber \\
		&b)\mkern+6mu\mathsf{L}f^*(\mathcal{G}_1^\bullet)\overset{\mathsf{L}}{\otimes}_{\mathcal{O}_Y}\!\mathsf{L}f^*(\mathcal{G}_2^\bullet)\cong \mathsf{L}f^*(\mathcal{G}_1^\bullet\overset{\mathsf{L}}{\otimes}_{\mathcal{O}_Y}\!\mathcal{G}_2^\bullet) &\big(X,Y\text{ projective}\big) \nonumber \\
		&c)\mkern+6mu\textup{Hom}_{\mathsf{D^b}(X)}(\mathsf{L}f^*(\mathcal{G}^\bullet),\mathcal{F}^\bullet)\cong \textup{Hom}_{\mathsf{D^b}(Y)}(\mathcal{G}^\bullet,\mathsf{R}f_*(\mathcal{F}^\bullet)) &\big(f\text{ projective}\big) \nonumber \\
		&d)\mkern+6mu\mathsf{R}\mathcal{H}om(\mathcal{F}_1^\bullet,\mathcal{F}_2^\bullet\overset{\mathsf{L}}{\otimes}_{\mathcal{O}_X}\!\mathcal{F}_3^\bullet)\cong\mathsf{R}\mathcal{H}om(\mathcal{F}_1^\bullet,\mathcal{F}_2^\bullet)\overset{\mathsf{L}}{\otimes}_{\mathcal{O}_X}\!\mathcal{F}_3^\bullet  &\big(X\!\in\!\textup{obj}(\mathsf{Sch}(\mathbb{K}))\text{ smooth} \nonumber \\
		& \mkern+187mu\cong\mathsf{R}\mathcal{H}om(\mathsf{R}\mathcal{H}om(\mathcal{F}_2^\bullet,\mathcal{F}_1^\bullet),\mathcal{F}_3^\bullet) &\text{and projective}\big) \nonumber \\
		&e)\mkern+6mu\mathsf{R}\mathcal{H}om(\mathcal{F}_1^\bullet\overset{\mathsf{L}}{\otimes}_{\mathcal{O}_X}\!\mathcal{F}_2^\bullet,\mathcal{F}_3^\bullet)\cong\mathsf{R}\mathcal{H}om(\mathcal{F}_1^\bullet,\mathsf{R}\mathcal{H}om(\mathcal{F}_2^\bullet,\mathcal{F}_3^\bullet)) & (\text{same}) \nonumber \\
		&f)\mkern+6mu(\mathcal{F}_1^\bullet)^\vee\overset{\mathsf{L}}{\otimes}_{\mathcal{O}_X}\!\mathcal{F}_2^\bullet\cong \mathsf{R}\mathcal{H}om(\mathcal{F}_1^\bullet,\mathcal{F}_2^\bullet) & (\text{same}) \nonumber \\
		&g)\mkern+6mu\mathcal{F}^\bullet\cong((\mathcal{F}^\bullet)^{\vee})^\vee\!\!\cong \mathsf{R}\mathcal{H}om(\mathsf{R}\mathcal{H}om(\mathcal{F}^\bullet,\mathcal{O}_X),\mathcal{O}_X) & (\text{same}) \nonumber \\
		&h)\mkern+6mu\mathsf{L}f^*(\mathsf{R}\mathcal{H}om_{\mathcal{O}_Y}\!(\mathcal{G}_1^\bullet,\mathcal{G}_2^\bullet))\!\cong\! \mathsf{R}\mathcal{H}om_{\mathcal{O}_X}\!(\mathsf{L}f^*(\mathcal{G}_1^\bullet),\mathsf{L}f^*(\mathcal{G}_2^\bullet)) & \big(X,Y\text{ proj.},\,\mathcal{G}_1^\bullet\in\mathsf{D^{\!-}}\!(Y)\big) \nonumber \\
		&i)\mkern+6mup_{X*}(p_Y^*\mathcal{G}^\bullet)\cong\mathsf{R}\Gamma(Y,\mathcal{G}^\bullet)\otimes_{\mathcal{O}_X}\!\mathcal{O}_X & \big(X\!\xleftarrow{p_X}\!X\!\times \!Y\!\xrightarrow{p_Y}\!Y\text{ proj.}\big) \nonumber \\
		&j)\mkern+6mu\mathsf{R}\Gamma(X\!\times\! Y,p_X^*\mathcal{F}\overset{\mathsf{L}}{\otimes}_{\mathcal{O}_X}\!p_Y^*\mathcal{G}^\bullet)\cong\mathsf{R}\Gamma(X,\mathcal{F}^\bullet)\otimes\mathsf{R}\Gamma(Y,\mathcal{G}^\bullet) & (\text{same})  \nonumber
	\end{align}\endgroup\end{small}
\!\!\!举其要者：a\textup) 是\textbf{投影公式}\index{projection formula}，c\textup) 是\textbf{伴随公式}\index{adjunction formula} \textup(实际上断言 $\mathsf{L}f_*$ 与 $\mathsf{R}f^*$ 是伴随函子\textup)，i\textup) 是\textbf{平坦基变换}\index{flat base change} 的特殊情形 \textup(见 \textup{\cite[equation (3.18)]{[Huy06]}}\textup)，而 j\textup) 诱导出\textbf{层上同调的 Künneth 公式}\index{Künneth formula for sheaf cohomology}。
\end{Pro} 

% [Tutor Note: 这份清单是 B-侧导出范畴的"演算手册"，但每条都有隐藏前提：a) 用投影公式需 $f$ 固有（否则 $\mathsf{R}f_*$ 不保持凝聚性）；d)、f)、g) 需要相应的复形是完美的（perfect，即局部同构于有限秩局部自由复形），否则等式不成立；g) 的对偶复原即"完美复形"的刻画。作者提示的证明策略是把复形逐项换成局部自由层（适配类）后逐茎验证。]
